\rSec1[expected]{Expected objects}
\indexlibraryglobal{expected}%

\rSec2[expected.general]{General}

\pnum
Subclause \ref{expected} describes the class template \tcode{expected}
that represents expected objects.
An \tcode{expected<T, E>} object holds
an object of type \tcode{T} or an object of type \tcode{E} and
manages the lifetime of the contained objects.

\rSec2[expected.syn]{Header \tcode{<expected>} synopsis}

\indexheader{expected}%
\indexlibraryglobal{unexpect_t}%
\indexlibraryglobal{unexpect}%
\begin{codeblock}
// mostly freestanding
namespace std {
  // \ref{expected.unexpected}, class template \tcode{unexpected}
  template<class E> class unexpected;

  // \ref{expected.bad}, class template \tcode{bad_expected_access}
  template<class E> class bad_expected_access;

  // \ref{expected.bad.void}, specialization for \tcode{void}
  template<> class bad_expected_access<void>;

  // in-place construction of unexpected values
  struct unexpect_t {
    explicit unexpect_t() = default;
  };
  inline constexpr unexpect_t unexpect{};

  // \ref{expected.expected}, class template \tcode{expected}
  template<class T, class E> class expected;                                // partially freestanding

  // \ref{expected.void}, partial specialization of \tcode{expected} for \tcode{void} types
  template<class T, class E> requires is_void_v<T> class expected<T, E>;    // partially freestanding
}
\end{codeblock}

\rSec2[expected.unexpected]{Class template \tcode{unexpected}}

\rSec3[expected.un.general]{General}

\pnum
Subclause \ref{expected.unexpected} describes the class template \tcode{unexpected}
that represents unexpected objects stored in \tcode{expected} objects.

\indexlibraryglobal{unexpected}%
\begin{codeblock}
namespace std {
  template<class E>
  class unexpected {
  public:
    // \ref{expected.un.cons}, constructors
    constexpr unexpected(const unexpected&) = default;
    constexpr unexpected(unexpected&&) = default;
    template<class Err = E>
      constexpr explicit unexpected(Err&&);
    template<class... Args>
      constexpr explicit unexpected(in_place_t, Args&&...);
    template<class U, class... Args>
      constexpr explicit unexpected(in_place_t, initializer_list<U>, Args&&...);

    constexpr unexpected& operator=(const unexpected&) = default;
    constexpr unexpected& operator=(unexpected&&) = default;

    constexpr const E& error() const & noexcept;
    constexpr E& error() & noexcept;
    constexpr const E&& error() const && noexcept;
    constexpr E&& error() && noexcept;

    constexpr void swap(unexpected& other) noexcept(@\seebelow@);

    template<class E2>
      friend constexpr bool operator==(const unexpected&, const unexpected<E2>&);

    friend constexpr void swap(unexpected& x, unexpected& y) noexcept(noexcept(x.swap(y)));

  private:
    E @\exposidnc{unex}@;             // \expos
  };

  template<class E> unexpected(E) -> unexpected<E>;
}
\end{codeblock}

\pnum
A program that instantiates the definition of \tcode{unexpected} for
a non-object type\added{ other than an lvalue reference type},
an array type,
a specialization of \tcode{unexpected}, or
a cv-qualified type
is ill-formed.

\rSec3[expected.un.cons]{Constructors}

\indexlibraryctor{unexpected}%
\begin{itemdecl}
template<class Err = E>
  constexpr explicit unexpected(Err&& e);
\end{itemdecl}

\begin{itemdescr}
\pnum
\constraints
\begin{itemize}
\item
\tcode{is_same_v<remove_cvref_t<Err>, unexpected>} is \tcode{false}; and
\item
\tcode{is_same_v<remove_cvref_t<Err>, in_place_t>} is \tcode{false}; and
\item
\tcode{is_constructible_v<E, Err>} is \tcode{true}.
\end{itemize}

\pnum
\effects
Direct-non-list-initializes \exposid{unex} with \tcode{std::forward<Err>(e)}.

\pnum
\throws
Any exception thrown by the initialization of \exposid{unex}.
\end{itemdescr}

\indexlibraryctor{unexpected}%
\begin{itemdecl}
template<class... Args>
  constexpr explicit unexpected(in_place_t, Args&&... args);
\end{itemdecl}

\begin{itemdescr}
\pnum
\constraints
\tcode{is_constructible_v<E, Args...>} is \tcode{true}.

\pnum
\effects
Direct-non-list-initializes
\exposid{unex} with \tcode{std::forward<Args>(args)...}.

\pnum
\throws
Any exception thrown by the initialization of \exposid{unex}.
\end{itemdescr}

\indexlibraryctor{unexpected}%
\begin{itemdecl}
template<class U, class... Args>
  constexpr explicit unexpected(in_place_t, initializer_list<U> il, Args&&... args);
\end{itemdecl}

\begin{itemdescr}
\pnum
\constraints
\tcode{is_constructible_v<E, initializer_list<U>\&, Args...>} is \tcode{true}.

\pnum
\effects
Direct-non-list-initializes
\exposid{unex} with \tcode{il, std::forward<Args>(args)...}.

\pnum
\throws
Any exception thrown by the initialization of \exposid{unex}.
\end{itemdescr}

\rSec3[expected.un.obs]{Observers}

\indexlibrarymember{error}{unexpected}%
\begin{itemdecl}
constexpr const E& error() const & noexcept;
constexpr E& error() & noexcept;
\end{itemdecl}

\begin{itemdescr}
\pnum
\returns
\exposid{unex}.
\end{itemdescr}

\indexlibrarymember{error}{unexpected}%
\begin{itemdecl}
constexpr E&& error() && noexcept;
constexpr const E&& error() const && noexcept;
\end{itemdecl}

\begin{itemdescr}
\pnum
\returns
\tcode{std::move(\exposid{unex})}.
\end{itemdescr}

\rSec3[expected.un.swap]{Swap}

\indexlibrarymember{swap}{unexpected}%
\begin{itemdecl}
constexpr void swap(unexpected& other) noexcept(is_nothrow_swappable_v<E>);
\end{itemdecl}

\begin{itemdescr}
\pnum
\mandates
\tcode{is_swappable_v<E>} is \tcode{true}.

\pnum
\effects
Equivalent to:
\tcode{using std::swap; swap(\exposid{unex}, other.\exposid{unex});}
\end{itemdescr}

\begin{itemdecl}
friend constexpr void swap(unexpected& x, unexpected& y) noexcept(noexcept(x.swap(y)));
\end{itemdecl}

\begin{itemdescr}
\pnum
\constraints
\tcode{is_swappable_v<E>} is \tcode{true}.

\pnum
\effects
Equivalent to \tcode{x.swap(y)}.
\end{itemdescr}

\rSec3[expected.un.eq]{Equality operator}

\indexlibrarymember{operator==}{unexpected}%
\begin{itemdecl}
template<class E2>
  friend constexpr bool operator==(const unexpected& x, const unexpected<E2>& y);
\end{itemdecl}

\begin{itemdescr}
\pnum
\mandates
The expression \tcode{x.error() == y.error()} is well-formed and
its result is convertible to \tcode{bool}.

\pnum
\returns
\tcode{x.error() == y.error()}.
\end{itemdescr}

\begin{addedblock}
\rSec3[expected.un.ref]{Partial specialization \tcode{unexpected<E\&>}}

\indexlibraryglobal{unexpected}%
\begin{codeblock}
template<class E>
class unexpected<E&> {
public:
  constexpr unexpected(const unexpected&) = default;
  constexpr unexpected(unexpected&&) = default;
  template<class G = E>
    constexpr explicit unexpected(G&&) noexcept;
  template<class G = E>
    constexpr explicit unexpected(in_place_t, G&&) noexcept;

  constexpr unexpected& operator=(const unexpected&) = default;
  constexpr unexpected& operator=(unexpected&&) = default;

  constexpr E& error() const noexcept;

  constexpr void swap(unexpected& other) noexcept;

  template<class E2>
    friend constexpr bool operator==(const unexpected&, const unexpected<E2>&);

  friend constexpr void swap(unexpected& x, unexpected& y) noexcept;

private:
  E* @\exposidnc{unex}@;            // \expos
};
\end{codeblock}

\pnum
An object of type \tcode{unexpected<E\&>} holds a pointer to an object of type
\tcode{E}. The referenced object is not owned by the \tcode{unexpected<E\&>}
object.

\pnum
A program that instantiates the definition of \tcode{unexpected<E\&>} for
an array type or a specialization of \tcode{unexpected} is ill-formed.
Unlike the primary template, \tcode{E} may be a cv-qualified type.

\indexlibraryctor{unexpected}%
\begin{itemdecl}
template<class G = E>
  constexpr explicit unexpected(G&& e) noexcept;
\end{itemdecl}

\begin{itemdescr}
\pnum
\constraints
\begin{itemize}
\item \tcode{is_same_v<remove_cvref_t<G>, unexpected>} is \tcode{false},
\item \tcode{is_same_v<remove_cvref_t<G>, in_place_t>} is \tcode{false},
\item \tcode{is_constructible_v<E\&, G>} is \tcode{true}, and
\item \tcode{reference_constructs_from_temporary_v<E\&, G>} is \tcode{false}.
\end{itemize}

\pnum
\effects
Initializes \exposid{unex} with
\tcode{addressof(static_cast<E\&>(std::forward<G>(e)))}.

\pnum
\remarks
A constructor for which
\tcode{reference_constructs_from_temporary_v<E\&, G>} is \tcode{true} ---
one that would bind \tcode{E\&} to a temporary --- is defined as deleted.
\end{itemdescr}

\indexlibraryctor{unexpected}%
\begin{itemdecl}
template<class G = E>
  constexpr explicit unexpected(in_place_t, G&& e) noexcept;
\end{itemdecl}

\begin{itemdescr}
\pnum
\constraints
\tcode{is_constructible_v<E\&, G>} is \tcode{true} and
\tcode{reference_constructs_from_temporary_v<E\&, G>} is \tcode{false}.

\pnum
\effects
Initializes \exposid{unex} with
\tcode{addressof(static_cast<E\&>(std::forward<G>(e)))}.
\end{itemdescr}

\indexlibrarymember{error}{unexpected}%
\begin{itemdecl}
constexpr E& error() const noexcept;
\end{itemdecl}

\begin{itemdescr}
\pnum
\returns
\tcode{*\exposid{unex}}.

\pnum
\remarks
The reference returned is not \tcode{const}-qualified even when
\tcode{*this} is \tcode{const}; the constness of the \tcode{unexpected<E\&>}
object does not propagate to the referenced object.
\end{itemdescr}

\indexlibrarymember{swap}{unexpected}%
\begin{itemdecl}
constexpr void swap(unexpected& other) noexcept;
\end{itemdecl}

\begin{itemdescr}
\pnum
\effects
Exchanges \exposid{unex} and \tcode{other.\exposid{unex}}.
\end{itemdescr}

\indexlibrarymember{operator==}{unexpected}%
\begin{itemdecl}
template<class E2>
  friend constexpr bool operator==(const unexpected& x, const unexpected<E2>& y);
\end{itemdecl}

\begin{itemdescr}
\pnum
\mandates
The expression \tcode{x.error() == y.error()} is well-formed and
its result is convertible to \tcode{bool}.

\pnum
\returns
\tcode{x.error() == y.error()}.
\end{itemdescr}

\indexlibrarymember{swap}{unexpected}%
\begin{itemdecl}
friend constexpr void swap(unexpected& x, unexpected& y) noexcept;
\end{itemdecl}

\begin{itemdescr}
\pnum
\effects
Equivalent to \tcode{x.swap(y)}.
\end{itemdescr}
\end{addedblock}

\rSec2[expected.bad]{Class template \tcode{bad_expected_access}}

\indexlibraryglobal{bad_expected_access}%
\begin{codeblock}
namespace std {
  template<class E>
  class bad_expected_access : public bad_expected_access<void> {
  public:
    constexpr explicit bad_expected_access(E);
    constexpr const char* what() const noexcept override;
    constexpr E& error() & noexcept;
    constexpr const E& error() const & noexcept;
    constexpr E&& error() && noexcept;
    constexpr const E&& error() const && noexcept;

  private:
    E @\exposidnc{unex}@;             // \expos
  };
}
\end{codeblock}

\pnum
The class template \tcode{bad_expected_access}
defines the type of objects thrown as exceptions to report the situation
where an attempt is made to access the value of an \tcode{expected<T, E>} object
for which \tcode{has_value()} is \tcode{false}.

\indexlibraryctor{bad_expected_access}%
\begin{itemdecl}
constexpr explicit bad_expected_access(E e);
\end{itemdecl}

\begin{itemdescr}
\pnum
\effects
Initializes \exposid{unex} with \tcode{std::move(e)}.
\end{itemdescr}

\indexlibrarymember{error}{bad_expected_access}%
\begin{itemdecl}
constexpr const E& error() const & noexcept;
constexpr E& error() & noexcept;
\end{itemdecl}

\begin{itemdescr}
\pnum
\returns
\exposid{unex}.
\end{itemdescr}

\indexlibrarymember{error}{bad_expected_access}%
\begin{itemdecl}
constexpr E&& error() && noexcept;
constexpr const E&& error() const && noexcept;
\end{itemdecl}

\begin{itemdescr}
\pnum
\returns
\tcode{std::move(\exposid{unex})}.
\end{itemdescr}

\indexlibrarymember{what}{bad_expected_access}%
\begin{itemdecl}
constexpr const char* what() const noexcept override;
\end{itemdecl}

\begin{itemdescr}
\pnum
\returns
An \impldef{return value of \tcode{bad_expected_access<E>::what}} \ntbs,
which during constant evaluation is encoded with
the ordinary literal encoding\iref{lex.ccon}.
\end{itemdescr}

\rSec2[expected.bad.void]{Class template specialization \tcode{bad_expected_access<void>}}

\begin{codeblock}
namespace std {
  template<>
  class bad_expected_access<void> : public exception {
  protected:
    constexpr bad_expected_access() noexcept;
    constexpr bad_expected_access(const bad_expected_access&) noexcept;
    constexpr bad_expected_access(bad_expected_access&&) noexcept;
    constexpr bad_expected_access& operator=(const bad_expected_access&) noexcept;
    constexpr bad_expected_access& operator=(bad_expected_access&&) noexcept;
    constexpr ~bad_expected_access();

  public:
    constexpr const char* what() const noexcept override;
  };
}
\end{codeblock}

\indexlibrarymember{what}{bad_expected_access}%
\begin{itemdecl}
constexpr const char* what() const noexcept override;
\end{itemdecl}

\begin{itemdescr}
\pnum
\returns
An \impldef{return value of \tcode{bad_expected_access<void>::what}} \ntbs,
which during constant evaluation is encoded with
the ordinary literal encoding\iref{lex.ccon}.
\end{itemdescr}

\rSec2[expected.expected]{Class template \tcode{expected}}

\rSec3[expected.object.general]{General}

\begin{codeblock}
namespace std {
  template<class T, class E>
  class expected {
  public:
    using @\libmember{value_type}{expected}@ = T;
    using @\libmember{error_type}{expected}@ = E;
    using @\libmember{unexpected_type}{expected}@ = unexpected<E>;

    template<class U>
    using @\libmember{rebind}{expected}@ = expected<U, error_type>;

    // \ref{expected.object.cons}, constructors
    constexpr expected();
    constexpr expected(const expected&);
    constexpr expected(expected&&) noexcept(@\seebelow@);
    template<class U, class G>
      constexpr explicit(@\seebelow@) expected(const expected<U, G>&);
    template<class U, class G>
      constexpr explicit(@\seebelow@) expected(expected<U, G>&&);

    template<class U = remove_cv_t<T>>
      constexpr explicit(@\seebelow@) expected(U&& v);

    template<class G>
      constexpr explicit(@\seebelow@) expected(const unexpected<G>&);
    template<class G>
      constexpr explicit(@\seebelow@) expected(unexpected<G>&&);

    template<class... Args>
      constexpr explicit expected(in_place_t, Args&&...);
    template<class U, class... Args>
      constexpr explicit expected(in_place_t, initializer_list<U>, Args&&...);
    template<class... Args>
      constexpr explicit expected(unexpect_t, Args&&...);
    template<class U, class... Args>
      constexpr explicit expected(unexpect_t, initializer_list<U>, Args&&...);

    // \ref{expected.object.dtor}, destructor
    constexpr ~expected();

    // \ref{expected.object.assign}, assignment
    constexpr expected& operator=(const expected&);
    constexpr expected& operator=(expected&&) noexcept(@\seebelow@);
    template<class U = remove_cv_t<T>> constexpr expected& operator=(U&&);
    template<class G>
      constexpr expected& operator=(const unexpected<G>&);
    template<class G>
      constexpr expected& operator=(unexpected<G>&&);

    template<class... Args>
      constexpr T& emplace(Args&&...) noexcept;
    template<class U, class... Args>
      constexpr T& emplace(initializer_list<U>, Args&&...) noexcept;

    // \ref{expected.object.swap}, swap
    constexpr void swap(expected&) noexcept(@\seebelow@);
    friend constexpr void swap(expected& x, expected& y) noexcept(noexcept(x.swap(y)));

    // \ref{expected.object.obs}, observers
    constexpr const T* operator->() const noexcept;
    constexpr T* operator->() noexcept;
    constexpr const T& operator*() const & noexcept;
    constexpr T& operator*() & noexcept;
    constexpr const T&& operator*() const && noexcept;
    constexpr T&& operator*() && noexcept;
    constexpr explicit operator bool() const noexcept;
    constexpr bool has_value() const noexcept;
    constexpr const T& value() const &;                                     // freestanding-deleted
    constexpr T& value() &;                                                 // freestanding-deleted
    constexpr const T&& value() const &&;                                   // freestanding-deleted
    constexpr T&& value() &&;                                               // freestanding-deleted
    constexpr const E& error() const & noexcept;
    constexpr E& error() & noexcept;
    constexpr const E&& error() const && noexcept;
    constexpr E&& error() && noexcept;
    template<class U = remove_cv_t<T>> constexpr T value_or(U&&) const &;
    template<class U = remove_cv_t<T>> constexpr T value_or(U&&) &&;
    template<class G = E> constexpr E error_or(G&&) const &;
    template<class G = E> constexpr E error_or(G&&) &&;

    // \ref{expected.object.monadic}, monadic operations
    template<class F> constexpr auto and_then(F&& f) &;
    template<class F> constexpr auto and_then(F&& f) &&;
    template<class F> constexpr auto and_then(F&& f) const &;
    template<class F> constexpr auto and_then(F&& f) const &&;
    template<class F> constexpr auto or_else(F&& f) &;
    template<class F> constexpr auto or_else(F&& f) &&;
    template<class F> constexpr auto or_else(F&& f) const &;
    template<class F> constexpr auto or_else(F&& f) const &&;
    template<class F> constexpr auto transform(F&& f) &;
    template<class F> constexpr auto transform(F&& f) &&;
    template<class F> constexpr auto transform(F&& f) const &;
    template<class F> constexpr auto transform(F&& f) const &&;
    template<class F> constexpr auto transform_error(F&& f) &;
    template<class F> constexpr auto transform_error(F&& f) &&;
    template<class F> constexpr auto transform_error(F&& f) const &;
    template<class F> constexpr auto transform_error(F&& f) const &&;

    // \ref{expected.object.eq}, equality operators
    template<class T2, class E2> requires (!is_void_v<T2>)
      friend constexpr bool operator==(const expected& x, const expected<T2, E2>& y);
    template<class T2>
      friend constexpr bool operator==(const expected&, const T2&);
    template<class E2>
      friend constexpr bool operator==(const expected&, const unexpected<E2>&);

  private:
    bool @\exposid{has_val}@;           // \expos
    union {
      remove_cv_t<T> @\exposid{val}@;   // \expos
      unexpected<E> @\exposid{unex}@;   // \expos
    };
  };
}
\end{codeblock}

\pnum
Any object of type \tcode{expected<T, E>} either
contains a value of type \tcode{T} or
a value of type \tcode{E}
nested within\iref{intro.object} it.
Member \exposid{has_val} indicates whether the \tcode{expected<T, E>} object
contains an object of type \tcode{T}.

\begin{addedblock}
\pnum
When \tcode{has_value()} is \tcode{false}, the error is
\tcode{\exposid{unex}.error()}. The error is held as an \tcode{unexpected<E>},
and not as an \tcode{E}, so that \tcode{E} may be an lvalue reference type: an
\tcode{E\&} cannot be a union member, whereas \tcode{unexpected<E\&>} holds a
pointer to an external object\iref{expected.un.ref}.
\end{addedblock}

\pnum
A type \tcode{T} is a \term{valid value type for \tcode{expected}},
if \tcode{remove_cv_t<T>} is \tcode{void}
or a complete non-array object type that is not \tcode{in_place_t},
\tcode{unexpect_t},
or a specialization of \tcode{unexpected}.
A program which instantiates class template \tcode{expected<T, E>}
with an argument \tcode{T} that is not a valid value
type for \tcode{expected} is ill-formed.
A program that instantiates
the definition of the template \tcode{expected<T, E>}
with a type for the \tcode{E} parameter
that is not a valid template argument for \tcode{unexpected} is ill-formed.

\pnum
When \tcode{T} is not \cv{} \tcode{void}, it shall meet
the \oldconcept{Destructible} requirements (\tref{cpp17.destructible}).
\added{If \tcode{E} is not an lvalue reference type, }\tcode{E} shall meet
the \oldconcept{Destructible} requirements.

\rSec3[expected.object.cons]{Constructors}

\pnum
The exposition-only variable template \exposid{converts-from-any-cvref}
defined in \ref{optional.ctor}
is used by some constructors for \tcode{expected}.

\indexlibraryctor{expected}%
\begin{itemdecl}
constexpr expected();
\end{itemdecl}

\begin{itemdescr}
\pnum
\constraints
\tcode{is_default_constructible_v<T>} is \tcode{true}.

\pnum
\effects
Value-initializes \exposid{val}.

\pnum
\ensures
\tcode{has_value()} is \tcode{true}.

\pnum
\throws
Any exception thrown by the initialization of \exposid{val}.
\end{itemdescr}

\indexlibraryctor{expected}%
\begin{itemdecl}
constexpr expected(const expected& rhs);
\end{itemdecl}

\begin{itemdescr}
\pnum
\effects
If \tcode{rhs.has_value()} is \tcode{true},
direct-non-list-initializes \exposid{val} with \tcode{*rhs}.
Otherwise, direct-non-list-initializes \exposid{unex} with \tcode{rhs.error()}.

\pnum
\ensures
\tcode{rhs.has_value() == this->has_value()}.

\pnum
\throws
Any exception thrown by the initialization of \exposid{val} or \exposid{unex}.

\pnum
\remarks
This constructor is defined as deleted unless
\begin{itemize}
\item
\tcode{is_copy_constructible_v<T>} is \tcode{true} and
\item
\tcode{is_copy_constructible_v<E>} is \tcode{true}.
\end{itemize}

\pnum
This constructor is trivial if
\begin{itemize}
\item
\tcode{is_trivially_copy_constructible_v<T>} is \tcode{true} and
\item
\tcode{is_trivially_copy_constructible_v<E>} is \tcode{true}.
\end{itemize}
\end{itemdescr}

\indexlibraryctor{expected}%
\begin{itemdecl}
constexpr expected(expected&& rhs) noexcept(@\seebelow@);
\end{itemdecl}

\begin{itemdescr}
\pnum
\constraints
\begin{itemize}
\item
\tcode{is_move_constructible_v<T>} is \tcode{true} and
\item
\tcode{is_move_constructible_v<E>} is \tcode{true}.
\end{itemize}

\pnum
\effects
If \tcode{rhs.has_value()} is \tcode{true},
direct-non-list-initializes \exposid{val} with \tcode{std::move(*rhs)}.
Otherwise,
direct-non-list-initializes \exposid{unex} with \tcode{std::move(rhs.error())}.

\pnum
\ensures
\tcode{rhs.has_value()} is unchanged;
\tcode{rhs.has_value() == this->has_value()} is \tcode{true}.

\pnum
\throws
Any exception thrown by the initialization of \exposid{val} or \exposid{unex}.

\pnum
\remarks
The exception specification is equivalent to
\tcode{is_nothrow_move_constructible_v<T> \&\&
is_nothrow_move_constructible_v<E>}.

\pnum
This constructor is trivial if
\begin{itemize}
\item
\tcode{is_trivially_move_constructible_v<T>} is \tcode{true} and
\item
\tcode{is_trivially_move_constructible_v<E>} is \tcode{true}.
\end{itemize}
\end{itemdescr}

\indexlibraryctor{expected}%
\begin{itemdecl}
template<class U, class G>
  constexpr explicit(@\seebelow@) expected(const expected<U, G>& rhs);
template<class U, class G>
  constexpr explicit(@\seebelow@) expected(expected<U, G>&& rhs);
\end{itemdecl}

\begin{itemdescr}
\pnum
Let:
\begin{itemize}
\item
\tcode{UF} be \tcode{const U\&} for the first overload and
\tcode{U} for the second overload.
\item
\tcode{GF} be \tcode{const G\&} for the first overload and
\tcode{G} for the second overload.
\end{itemize}

\pnum
\constraints
\begin{itemize}
\item
\tcode{is_constructible_v<T, UF>} is \tcode{true}; and
\item
\tcode{is_constructible_v<E, GF>} is \tcode{true}; and
\item
if \tcode{T} is not \cv{} \tcode{bool},
\tcode{\exposid{converts-from-any-cvref}<T, expected<U, G>>} is \tcode{false}; and
\item
\tcode{is_constructible_v<unexpected<E>, expected<U, G>\&>} is \tcode{false}; and
\item
\tcode{is_constructible_v<unexpected<E>, expected<U, G>>} is \tcode{false}; and
\item
\tcode{is_constructible_v<unexpected<E>, const expected<U, G>\&>} is \tcode{false}; and
\item
\tcode{is_constructible_v<unexpected<E>, const expected<U, G>>} is \tcode{false}.
\end{itemize}

\pnum
\effects
If \tcode{rhs.has_value()},
direct-non-list-initializes \exposid{val} with \tcode{std::forward<UF>(*rhs)}.
Otherwise,
direct-non-list-initializes \exposid{unex} with \tcode{std::forward<GF>(rhs.error())}.

\pnum
\ensures
\tcode{rhs.has_value()} is unchanged;
\tcode{rhs.has_value() == this->has_value()} is \tcode{true}.

\pnum
\throws
Any exception thrown by the initialization of \exposid{val} or \exposid{unex}.

\pnum
\remarks
The expression inside \tcode{explicit} is equivalent to
\tcode{!is_convertible_v<UF, T> || !is_convertible_v<GF, E>}.
\end{itemdescr}

\indexlibraryctor{expected}%
\begin{itemdecl}
template<class U = remove_cv_t<T>>
  constexpr explicit(!is_convertible_v<U, T>) expected(U&& v);
\end{itemdecl}

\begin{itemdescr}
\pnum
\constraints
\begin{itemize}
\item
\tcode{is_same_v<remove_cvref_t<U>, in_place_t>} is \tcode{false}; and
\item
\tcode{is_same_v<remove_cvref_t<U>, expected>} is \tcode{false}; and
\item
\tcode{is_same_v<remove_cvref_t<U>, unexpect_t>} is \tcode{false}; and
\item
\tcode{remove_cvref_t<U>} is not a specialization of \tcode{unexpected}; and
\item
\tcode{is_constructible_v<T, U>} is \tcode{true}; and
\item
if \tcode{T} is \cv{} \tcode{bool},
\tcode{remove_cvref_t<U>} is not a specialization of \tcode{expected}.
\end{itemize}

\pnum
\effects
Direct-non-list-initializes \exposid{val} with \tcode{std::forward<U>(v)}.

\pnum
\ensures
\tcode{has_value()} is \tcode{true}.

\pnum
\throws
Any exception thrown by the initialization of \exposid{val}.
\end{itemdescr}

\indexlibraryctor{expected}%
\begin{itemdecl}
template<class G>
  constexpr explicit(!is_convertible_v<const G&, E>) expected(const unexpected<G>& e);
template<class G>
  constexpr explicit(!is_convertible_v<G, E>) expected(unexpected<G>&& e);
\end{itemdecl}

\begin{itemdescr}
\pnum
Let \tcode{GF} be \tcode{const G\&} for the first overload and
\tcode{G} for the second overload.

\pnum
\constraints
\tcode{is_constructible_v<E, GF>} is \tcode{true}.

\pnum
\effects
Direct-non-list-initializes \exposid{unex} with \tcode{std::forward<GF>(e.error())}.

\pnum
\ensures
\tcode{has_value()} is \tcode{false}.

\pnum
\throws
Any exception thrown by the initialization of \exposid{unex}.
\end{itemdescr}

\indexlibraryctor{expected}%
\begin{itemdecl}
template<class... Args>
  constexpr explicit expected(in_place_t, Args&&... args);
\end{itemdecl}

\begin{itemdescr}
\pnum
\constraints
\tcode{is_constructible_v<T, Args...>} is \tcode{true}.

\pnum
\effects
Direct-non-list-initializes \exposid{val} with \tcode{std::forward<Args>(args)...}.

\pnum
\ensures
\tcode{has_value()} is \tcode{true}.

\pnum
\throws
Any exception thrown by the initialization of \exposid{val}.
\end{itemdescr}

\indexlibraryctor{expected}%
\begin{itemdecl}
template<class U, class... Args>
  constexpr explicit expected(in_place_t, initializer_list<U> il, Args&&... args);
\end{itemdecl}

\begin{itemdescr}
\pnum
\constraints
\tcode{is_constructible_v<T, initializer_list<U>\&, Args...>} is \tcode{true}.

\pnum
\effects
Direct-non-list-initializes \exposid{val} with
\tcode{il, std::forward<Args>(args)...}.

\pnum
\ensures
\tcode{has_value()} is \tcode{true}.

\pnum
\throws
Any exception thrown by the initialization of \exposid{val}.
\end{itemdescr}

\indexlibraryctor{expected}%
\begin{itemdecl}
template<class... Args>
  constexpr explicit expected(unexpect_t, Args&&... args);
\end{itemdecl}

\begin{itemdescr}
\pnum
\constraints
\tcode{is_constructible_v<E, Args...>} is \tcode{true}.

\pnum
\effects
Direct-non-list-initializes \exposid{unex} with
\tcode{std::forward<Args>(args)...}.

\pnum
\ensures
\tcode{has_value()} is \tcode{false}.

\pnum
\throws
Any exception thrown by the initialization of \exposid{unex}.
\end{itemdescr}

\indexlibraryctor{expected}%
\begin{itemdecl}
template<class U, class... Args>
  constexpr explicit expected(unexpect_t, initializer_list<U> il, Args&&... args);
\end{itemdecl}

\begin{itemdescr}
\pnum
\constraints
\tcode{is_constructible_v<E, initializer_list<U>\&, Args...>} is \tcode{true}.

\pnum
\effects
Direct-non-list-initializes \exposid{unex} with
\tcode{il, std::forward<Args>(args)...}.

\pnum
\ensures
\tcode{has_value()} is \tcode{false}.

\pnum
\throws
Any exception thrown by the initialization of \exposid{unex}.
\end{itemdescr}

\rSec3[expected.object.dtor]{Destructor}

\indexlibrarydtor{expected}%
\begin{itemdecl}
constexpr ~expected();
\end{itemdecl}

\begin{itemdescr}
\pnum
\effects
If \tcode{has_value()} is \tcode{true}, destroys \exposid{val},
otherwise destroys \exposid{unex}.

\pnum
\remarks
If \tcode{is_trivially_destructible_v<T>} is \tcode{true}, and
\tcode{is_trivially_destructible_v<E>} is \tcode{true},
then this destructor is a trivial destructor.
\end{itemdescr}

\rSec3[expected.object.assign]{Assignment}

\pnum
This subclause makes use of the following exposition-only function template:
\begin{codeblock}
template<class T, class U, class... Args>
constexpr void @\exposid{reinit-expected}@(T& newval, U& oldval, Args&&... args) {  // \expos
  if constexpr (is_nothrow_constructible_v<T, Args...>) {
    destroy_at(addressof(oldval));
    construct_at(addressof(newval), std::forward<Args>(args)...);
  } else if constexpr (is_nothrow_move_constructible_v<T>) {
    T tmp(std::forward<Args>(args)...);
    destroy_at(addressof(oldval));
    construct_at(addressof(newval), std::move(tmp));
  } else {
    U tmp(std::move(oldval));
    destroy_at(addressof(oldval));
    try {
      construct_at(addressof(newval), std::forward<Args>(args)...);
    } catch (...) {
      construct_at(addressof(oldval), std::move(tmp));
      throw;
    }
  }
}
\end{codeblock}

\indexlibrarymember{operator=}{expected}%
\begin{itemdecl}
constexpr expected& operator=(const expected& rhs);
\end{itemdecl}

\begin{itemdescr}
\pnum
\effects
\begin{itemize}
\item
If \tcode{this->has_value() \&\& rhs.has_value()} is \tcode{true},
equivalent to \tcode{\exposid{val} = *rhs}.
\item
Otherwise, if \tcode{this->has_value()} is \tcode{true}, equivalent to:
\begin{codeblock}
@\exposid{reinit-expected}@(@\exposid{unex}@, @\exposid{val}@, rhs.error())
\end{codeblock}
\item
Otherwise, if \tcode{rhs.has_value()} is \tcode{true}, equivalent to:
\begin{codeblock}
@\exposid{reinit-expected}@(@\exposid{val}@, @\exposid{unex}@, *rhs)
\end{codeblock}
\item
Otherwise, equivalent to \tcode{\exposid{unex} = rhs.\exposid{unex}}.
\end{itemize}
Then, if no exception was thrown,
equivalent to: \tcode{\exposid{has_val} = rhs.has_value(); return *this;}

\pnum
\returns
\tcode{*this}.

\pnum
\remarks
This operator is defined as deleted unless:
\begin{itemize}
\item
\tcode{is_copy_assignable_v<T>} is \tcode{true} and
\item
\tcode{is_copy_constructible_v<T>} is \tcode{true} and
\item
\tcode{is_copy_assignable_v<E>} is \tcode{true} and
\item
\tcode{is_copy_constructible_v<E>} is \tcode{true} and
\item
\tcode{is_nothrow_move_constructible_v<T> || is_nothrow_move_constructible_v<E>}
is \tcode{true}.
\end{itemize}

\pnum
This operator is trivial if:
\begin{itemize}
\item \tcode{is_trivially_copy_constructible_v<T>} is \tcode{true}, and
\item \tcode{is_trivially_copy_assignable_v<T>} is \tcode{true}, and
\item \tcode{is_trivially_destructible_v<T>} is \tcode{true}, and
\item \tcode{is_trivially_copy_constructible_v<E>} is \tcode{true}, and
\item \tcode{is_trivially_copy_assignable_v<E>} is \tcode{true}, and
\item \tcode{is_trivially_destructible_v<E>} is \tcode{true}.
\end{itemize}
\end{itemdescr}

\indexlibrarymember{operator=}{expected}%
\begin{itemdecl}
constexpr expected& operator=(expected&& rhs) noexcept(@\seebelow@);
\end{itemdecl}

\begin{itemdescr}
\pnum
\constraints
\begin{itemize}
\item
\tcode{is_move_constructible_v<T>} is \tcode{true} and
\item
\tcode{is_move_assignable_v<T>} is \tcode{true} and
\item
\tcode{is_move_constructible_v<E>} is \tcode{true} and
\item
\tcode{is_move_assignable_v<E>} is \tcode{true} and
\item
\tcode{is_nothrow_move_constructible_v<T> || is_nothrow_move_constructible_v<E>}
is \tcode{true}.
\end{itemize}

\pnum
\effects
\begin{itemize}
\item
If \tcode{this->has_value() \&\& rhs.has_value()} is \tcode{true},
equivalent to \tcode{\exposid{val} = std::move(*rhs)}.
\item
Otherwise, if \tcode{this->has_value()} is \tcode{true}, equivalent to:
\begin{codeblock}
@\exposid{reinit-expected}@(@\exposid{unex}@, @\exposid{val}@, std::move(rhs.error()))
\end{codeblock}
\item
Otherwise, if \tcode{rhs.has_value()} is \tcode{true}, equivalent to:
\begin{codeblock}
@\exposid{reinit-expected}@(@\exposid{val}@, @\exposid{unex}@, std::move(*rhs))
\end{codeblock}
\item
Otherwise, equivalent to \tcode{\exposid{unex} = std::move(rhs.\exposid{unex})}.
\end{itemize}
Then, if no exception was thrown,
equivalent to: \tcode{has_val = rhs.has_value(); return *this;}

\pnum
\returns
\tcode{*this}.

\pnum
\remarks
The exception specification is equivalent to:
\begin{codeblock}
is_nothrow_move_assignable_v<T> && is_nothrow_move_constructible_v<T> &&
is_nothrow_move_assignable_v<E> && is_nothrow_move_constructible_v<E>
\end{codeblock}

\pnum
This operator is trivial if:
\begin{itemize}
\item \tcode{is_trivially_move_constructible_v<T>} is \tcode{true}, and
\item \tcode{is_trivially_move_assignable_v<T>} is \tcode{true}, and
\item \tcode{is_trivially_destructible_v<T>} is \tcode{true}, and
\item \tcode{is_trivially_move_constructible_v<E>} is \tcode{true}, and
\item \tcode{is_trivially_move_assignable_v<E>} is \tcode{true}, and
\item \tcode{is_trivially_destructible_v<E>} is \tcode{true}.
\end{itemize}
\end{itemdescr}

\indexlibrarymember{operator=}{expected}%
\begin{itemdecl}
template<class U = remove_cv_t<T>>
  constexpr expected& operator=(U&& v);
\end{itemdecl}

\begin{itemdescr}
\pnum
\constraints
\begin{itemize}
\item
\tcode{is_same_v<expected, remove_cvref_t<U>>} is \tcode{false}; and
\item
\tcode{remove_cvref_t<U>} is not a specialization of \tcode{unexpected}; and
\item
\tcode{is_constructible_v<T, U>} is \tcode{true}; and
\item
\tcode{is_assignable_v<T\&, U>} is \tcode{true}; and
\item
\tcode{is_nothrow_constructible_v<T, U> || is_nothrow_move_constructible_v<T> ||\newline
is_nothrow_move_constructible_v<E>}
is \tcode{true}.
\end{itemize}

\pnum
\effects
\begin{itemize}
\item
If \tcode{has_value()} is \tcode{true},
equivalent to: \tcode{\exposid{val} = std::forward<U>(v);}
\item
Otherwise, equivalent to:
\begin{codeblock}
@\exposid{reinit-expected}@(@\exposid{val}@, @\exposid{unex}@, std::forward<U>(v));
@\exposid{has_val}@ = true;
\end{codeblock}
\end{itemize}

\pnum
\returns
\tcode{*this}.
\end{itemdescr}

\indexlibrarymember{operator=}{expected}%
\begin{itemdecl}
template<class G>
  constexpr expected& operator=(const unexpected<G>& e);
template<class G>
  constexpr expected& operator=(unexpected<G>&& e);
\end{itemdecl}

\begin{itemdescr}
\pnum
Let \tcode{GF} be \tcode{const G\&} for the first overload and
\tcode{G} for the second overload.

\pnum
\constraints
\begin{itemize}
\item
\tcode{is_constructible_v<E, GF>} is \tcode{true}; and
\item
\tcode{is_assignable_v<E\&, GF>} is \tcode{true}; and
\item
\tcode{is_nothrow_constructible_v<E, GF> || is_nothrow_move_constructible_v<T> ||\newline
is_nothrow_move_constructible_v<E>} is \tcode{true}.
\end{itemize}

\pnum
\effects
\begin{itemize}
\item
If \tcode{has_value()} is \tcode{true}, equivalent to:
\begin{codeblock}
@\exposid{reinit-expected}@(@\exposid{unex}@, @\exposid{val}@, std::forward<GF>(e.error()));
@\exposid{has_val}@ = false;
\end{codeblock}
\item
Otherwise, equivalent to:
\tcode{\exposid{unex} = unexpected<E>(std::forward<GF>(e.error()));}
\end{itemize}

\pnum
\returns
\tcode{*this}.
\end{itemdescr}

\indexlibrarymember{emplace}{expected}%
\begin{itemdecl}
template<class... Args>
  constexpr T& emplace(Args&&... args) noexcept;
\end{itemdecl}

\begin{itemdescr}
\pnum
\constraints
\tcode{is_nothrow_constructible_v<T, Args...>} is \tcode{true}.

\pnum
\effects
Equivalent to:
\begin{codeblock}
if (has_value()) {
  destroy_at(addressof(@\exposid{val}@));
} else {
  destroy_at(addressof(@\exposid{unex}@));
  @\exposid{has_val}@ = true;
}
return *construct_at(addressof(@\exposid{val}@), std::forward<Args>(args)...);
\end{codeblock}
\end{itemdescr}

\indexlibrarymember{emplace}{expected}%
\begin{itemdecl}
template<class U, class... Args>
  constexpr T& emplace(initializer_list<U> il, Args&&... args) noexcept;
\end{itemdecl}

\begin{itemdescr}
\pnum
\constraints
\tcode{is_nothrow_constructible_v<T, initializer_list<U>\&, Args...>}
is \tcode{true}.

\pnum
\effects
Equivalent to:
\begin{codeblock}
if (has_value()) {
  destroy_at(addressof(@\exposid{val}@));
} else {
  destroy_at(addressof(@\exposid{unex}@));
  @\exposid{has_val}@ = true;
}
return *construct_at(addressof(@\exposid{val}@), il, std::forward<Args>(args)...);
\end{codeblock}
\end{itemdescr}

\rSec3[expected.object.swap]{Swap}

\indexlibrarymember{swap}{expected}%
\begin{itemdecl}
constexpr void swap(expected& rhs) noexcept(@\seebelow@);
\end{itemdecl}

\begin{itemdescr}
\pnum
\constraints
\begin{itemize}
\item
\tcode{is_swappable_v<T>} is \tcode{true} and
\item
\tcode{is_swappable_v<E>} is \tcode{true} and
\item
\tcode{is_move_constructible_v<T> \&\& is_move_constructible_v<E>}
is \tcode{true}, and
\item
\tcode{is_nothrow_move_constructible_v<T> || is_nothrow_move_constructible_v<E>}
is \tcode{true}.
\end{itemize}

\pnum
\effects
See \tref{expected.object.swap}.

\begin{floattable}{\tcode{swap(expected\&)} effects}{expected.object.swap}
{lx{0.35\hsize}x{0.35\hsize}}
\topline
& \chdr{\tcode{this->has_value()}} & \rhdr{\tcode{!this->has_value()}} \\ \capsep
\lhdr{\tcode{rhs.has_value()}} &
  equivalent to: \tcode{using std::swap; swap(\exposid{val}, rhs.\exposid{val});} &
  calls \tcode{rhs.swap(*this)} \\
\lhdr{\tcode{!rhs.has_value()}} &
  \seebelow &
  equivalent to: \tcode{using std::swap; swap(\exposid{unex}, rhs.\exposid{unex});} \\
\end{floattable}

For the case where \tcode{rhs.has_value()} is \tcode{false} and
\tcode{this->has_value()} is \tcode{true}, equivalent to:
\begin{codeblock}
if constexpr (is_nothrow_move_constructible_v<unexpected<E>>) {
  unexpected<E> tmp(std::move(rhs.@\exposid{unex}@));
  destroy_at(addressof(rhs.@\exposid{unex}@));
  try {
    construct_at(addressof(rhs.@\exposid{val}@), std::move(@\exposid{val}@));
    destroy_at(addressof(@\exposid{val}@));
    construct_at(addressof(@\exposid{unex}@), std::move(tmp));
  } catch(...) {
    construct_at(addressof(rhs.@\exposid{unex}@), std::move(tmp));
    throw;
  }
} else {
  remove_cv_t<T> tmp(std::move(@\exposid{val}@));
  destroy_at(addressof(@\exposid{val}@));
  try {
    construct_at(addressof(@\exposid{unex}@), std::move(rhs.@\exposid{unex}@));
    destroy_at(addressof(rhs.@\exposid{unex}@));
    construct_at(addressof(rhs.@\exposid{val}@), std::move(tmp));
  } catch (...) {
    construct_at(addressof(@\exposid{val}@), std::move(tmp));
    throw;
  }
}
@\exposid{has_val}@ = false;
rhs.@\exposid{has_val}@ = true;
\end{codeblock}

\pnum
\throws
Any exception thrown by the expressions in the \Fundescx{Effects}.

\pnum
\remarks
The exception specification is equivalent to:
\begin{codeblock}
is_nothrow_move_constructible_v<T> && is_nothrow_swappable_v<T> &&
is_nothrow_move_constructible_v<unexpected<E>> && is_nothrow_swappable_v<unexpected<E>>
\end{codeblock}
\end{itemdescr}

\indexlibrarymember{swap}{expected}%
\begin{itemdecl}
friend constexpr void swap(expected& x, expected& y) noexcept(noexcept(x.swap(y)));
\end{itemdecl}

\begin{itemdescr}
\pnum
\effects
Equivalent to \tcode{x.swap(y)}.
\end{itemdescr}

\rSec3[expected.object.obs]{Observers}

\indexlibrarymember{operator->}{expected}%
\begin{itemdecl}
constexpr const T* operator->() const noexcept;
constexpr T* operator->() noexcept;
\end{itemdecl}

\begin{itemdescr}
\pnum
\hardexpects
\tcode{has_value()} is \tcode{true}.

\pnum
\returns
\tcode{addressof(\exposid{val})}.
\end{itemdescr}

\indexlibrarymember{operator*}{expected}%
\begin{itemdecl}
constexpr const T& operator*() const & noexcept;
constexpr T& operator*() & noexcept;
\end{itemdecl}

\begin{itemdescr}
\pnum
\hardexpects
\tcode{has_value()} is \tcode{true}.

\pnum
\returns
\exposid{val}.
\end{itemdescr}

\indexlibrarymember{operator*}{expected}%
\begin{itemdecl}
constexpr T&& operator*() && noexcept;
constexpr const T&& operator*() const && noexcept;
\end{itemdecl}

\begin{itemdescr}
\pnum
\hardexpects
\tcode{has_value()} is \tcode{true}.

\pnum
\returns
\tcode{std::move(\exposid{val})}.
\end{itemdescr}

\indexlibrarymember{operator bool}{expected}%
\indexlibrarymember{has_value}{expected}%
\begin{itemdecl}
constexpr explicit operator bool() const noexcept;
constexpr bool has_value() const noexcept;
\end{itemdecl}

\begin{itemdescr}
\pnum
\returns
\exposid{has_val}.
\end{itemdescr}

\indexlibrarymember{value}{expected}%
\begin{itemdecl}
constexpr const T& value() const &;
constexpr T& value() &;
\end{itemdecl}

\begin{itemdescr}
\pnum
\mandates
\tcode{is_copy_constructible_v<E>} is \tcode{true}.

\pnum
\returns
\exposid{val}, if \tcode{has_value()} is \tcode{true}.

\pnum
\throws
\tcode{bad_expected_access(as_const(error()))} if \tcode{has_value()} is \tcode{false}.
\end{itemdescr}

\indexlibrarymember{value}{expected}%
\begin{itemdecl}
constexpr T&& value() &&;
constexpr const T&& value() const &&;
\end{itemdecl}

\begin{itemdescr}
\pnum
\mandates
\tcode{is_copy_constructible_v<E>} is \tcode{true} and
\tcode{is_constructible_v<E, decltype(std::\linebreak{}move(error()))>} is \tcode{true}.

\pnum
\returns
\tcode{std::move(\exposid{val})}, if \tcode{has_value()} is \tcode{true}.

\pnum
\throws
\tcode{bad_expected_access(std::move(error()))}
if \tcode{has_value()} is \tcode{false}.
\end{itemdescr}

\indexlibrarymember{error}{expected}%
\begin{itemdecl}
constexpr const E& error() const & noexcept;
constexpr E& error() & noexcept;
\end{itemdecl}

\begin{itemdescr}
\pnum
\hardexpects
\tcode{has_value()} is \tcode{false}.

\pnum
\returns
\tcode{\exposid{unex}.error()}.
\end{itemdescr}

\indexlibrarymember{error}{expected}%
\begin{itemdecl}
constexpr E&& error() && noexcept;
constexpr const E&& error() const && noexcept;
\end{itemdecl}

\begin{itemdescr}
\pnum
\hardexpects
\tcode{has_value()} is \tcode{false}.

\pnum
\returns
\tcode{std::move(\exposid{unex}).error()}.
\end{itemdescr}

\indexlibrarymember{value_or}{expected}%
\begin{itemdecl}
template<class U = remove_cv_t<T>> constexpr T value_or(U&& v) const &;
\end{itemdecl}

\begin{itemdescr}
\pnum
\mandates
\tcode{is_copy_constructible_v<T>} is \tcode{true} and
\tcode{is_convertible_v<U, T>} is \tcode{true}.

\pnum
\returns
\tcode{has_value() ? **this : static_cast<T>(std::forward<U>(v))}.
\end{itemdescr}

\indexlibrarymember{value_or}{expected}%
\begin{itemdecl}
template<class U = remove_cv_t<T>> constexpr T value_or(U&& v) &&;
\end{itemdecl}

\begin{itemdescr}
\pnum
\mandates
\tcode{is_move_constructible_v<T>} is \tcode{true} and
\tcode{is_convertible_v<U, T>} is \tcode{true}.

\pnum
\returns
\tcode{has_value() ? std::move(**this) : static_cast<T>(std::forward<U>(v))}.
\end{itemdescr}

\indexlibrarymember{error_or}{expected}%
\begin{itemdecl}
template<class G = E> constexpr E error_or(G&& e) const &;
\end{itemdecl}

\begin{itemdescr}
\pnum
\mandates
\tcode{is_copy_constructible_v<E>} is \tcode{true} and
\tcode{is_convertible_v<G, E>} is \tcode{true}.

\pnum
\returns
\tcode{std::forward<G>(e)} if \tcode{has_value()} is \tcode{true},
\tcode{error()} otherwise.
\end{itemdescr}

\indexlibrarymember{error_or}{expected}%
\begin{itemdecl}
template<class G = E> constexpr E error_or(G&& e) &&;
\end{itemdecl}

\begin{itemdescr}
\pnum
\mandates
\tcode{is_move_constructible_v<E>} is \tcode{true} and
\tcode{is_convertible_v<G, E>} is \tcode{true}.

\pnum
\returns
\tcode{std::forward<G>(e)} if \tcode{has_value()} is \tcode{true},
\tcode{std::move(error())} otherwise.
\end{itemdescr}

\rSec3[expected.object.monadic]{Monadic operations}

\indexlibrarymember{and_then}{expected}%
\begin{itemdecl}
template<class F> constexpr auto and_then(F&& f) &;
template<class F> constexpr auto and_then(F&& f) const &;
\end{itemdecl}

\begin{itemdescr}
\pnum
Let \tcode{U} be \tcode{remove_cvref_t<invoke_result_t<F, decltype((\exposid{val}))>>}.

\pnum
\constraints
\tcode{is_constructible_v<E, decltype(error())>} is \tcode{true}.

\pnum
\mandates
\tcode{U} is a specialization of \tcode{expected} and
\tcode{is_same_v<typename U::error_type, E>} is \tcode{true}.

\pnum
\effects
Equivalent to:
\begin{codeblock}
if (has_value())
  return invoke(std::forward<F>(f), @\exposid{val}@);
else
  return U(unexpect, error());
\end{codeblock}
\end{itemdescr}

\indexlibrarymember{and_then}{expected}%
\begin{itemdecl}
template<class F> constexpr auto and_then(F&& f) &&;
template<class F> constexpr auto and_then(F&& f) const &&;
\end{itemdecl}

\begin{itemdescr}
\pnum
Let \tcode{U} be
\tcode{remove_cvref_t<invoke_result_t<F, decltype(std::move(\exposid{val}))>>}.

\pnum
\constraints
\tcode{is_constructible_v<E, decltype(std::move(error()))>} is \tcode{true}.

\pnum
\mandates
\tcode{U} is a specialization of \tcode{expected} and
\tcode{is_same_v<typename U::error_type, E>} is \tcode{true}.

\pnum
\effects
Equivalent to:
\begin{codeblock}
if (has_value())
  return invoke(std::forward<F>(f), std::move(@\exposid{val}@));
else
  return U(unexpect, std::move(error()));
\end{codeblock}
\end{itemdescr}

\indexlibrarymember{or_else}{expected}%
\begin{itemdecl}
template<class F> constexpr auto or_else(F&& f) &;
template<class F> constexpr auto or_else(F&& f) const &;
\end{itemdecl}

\begin{itemdescr}
\pnum
Let \tcode{G} be \tcode{remove_cvref_t<invoke_result_t<F, decltype(error())>>}.

\pnum
\constraints
\tcode{is_constructible_v<T, decltype((\exposid{val}))>} is \tcode{true}.

\pnum
\mandates
\tcode{G} is a specialization of \tcode{expected} and
\tcode{is_same_v<typename G::value_type, T>} is \tcode{true}.

\pnum
\effects
Equivalent to:
\begin{codeblock}
if (has_value())
  return G(in_place, @\exposid{val}@);
else
  return invoke(std::forward<F>(f), error());
\end{codeblock}
\end{itemdescr}

\indexlibrarymember{or_else}{expected}%
\begin{itemdecl}
template<class F> constexpr auto or_else(F&& f) &&;
template<class F> constexpr auto or_else(F&& f) const &&;
\end{itemdecl}

\begin{itemdescr}
\pnum
Let \tcode{G} be
\tcode{remove_cvref_t<invoke_result_t<F, decltype(std::move(error()))>>}.

\pnum
\constraints
\tcode{is_constructible_v<T, decltype(std::move(\exposid{val}))>} is \tcode{true}.

\pnum
\mandates
\tcode{G} is a specialization of \tcode{expected} and
\tcode{is_same_v<typename G::value_type, T>} is \tcode{true}.

\pnum
\effects
Equivalent to:
\begin{codeblock}
if (has_value())
  return G(in_place, std::move(@\exposid{val}@));
else
  return invoke(std::forward<F>(f), std::move(error()));
\end{codeblock}
\end{itemdescr}

\indexlibrarymember{transform}{expected}%
\begin{itemdecl}
template<class F> constexpr auto transform(F&& f) &;
template<class F> constexpr auto transform(F&& f) const &;
\end{itemdecl}

\begin{itemdescr}
\pnum
Let \tcode{U} be
\tcode{remove_cv_t<invoke_result_t<F, decltype((\exposid{val}))>>}.

\pnum
\constraints
\tcode{is_constructible_v<E, decltype(error())>} is \tcode{true}.

\pnum
\mandates
\tcode{U} is a valid value type for \tcode{expected}.
If \tcode{is_void_v<U>} is \tcode{false},
the declaration
\begin{codeblock}
U u(invoke(std::forward<F>(f), @\exposid{val}@));
\end{codeblock}
is well-formed.

\pnum
\effects
\begin{itemize}
\item
If \tcode{has_value()} is \tcode{false}, returns
\tcode{expected<U, E>(unexpect, error())}.
\item
Otherwise, if \tcode{is_void_v<U>} is \tcode{false}, returns an
\tcode{expected<U, E>} object whose \exposid{has_val} member is \tcode{true}
and \exposid{val} member is direct-non-list-initialized with
\tcode{invoke(std::forward<F>(f), \exposid{val})}.
\item
Otherwise, evaluates \tcode{invoke(std::forward<F>(f), \exposid{val})} and then
returns \tcode{expected<U, E>()}.
\end{itemize}
\end{itemdescr}

\indexlibrarymember{transform}{expected}%
\begin{itemdecl}
template<class F> constexpr auto transform(F&& f) &&;
template<class F> constexpr auto transform(F&& f) const &&;
\end{itemdecl}

\begin{itemdescr}
\pnum
Let \tcode{U} be
\tcode{remove_cv_t<invoke_result_t<F, decltype(std::move(\exposid{val}))>>}.

\pnum
\constraints
\tcode{is_constructible_v<E, decltype(std::move(error()))>} is \tcode{true}.

\pnum
\mandates
\tcode{U} is a valid value type for \tcode{expected}. If \tcode{is_void_v<U>} is
\tcode{false}, the declaration
\begin{codeblock}
U u(invoke(std::forward<F>(f), std::move(@\exposid{val}@)));
\end{codeblock}
is well-formed.

\pnum
\effects
\begin{itemize}
\item
If \tcode{has_value()} is \tcode{false}, returns
\tcode{expected<U, E>(unexpect, std::move(error()))}.
\item
Otherwise, if \tcode{is_void_v<U>} is \tcode{false}, returns an
\tcode{expected<U, E>} object whose \exposid{has_val} member is \tcode{true}
and \exposid{val} member is direct-non-list-initialized with
\tcode{invoke(std::forward<F>(f), std::move(\exposid{val}))}.
\item
Otherwise, evaluates \tcode{invoke(std::forward<F>(f), std::move(\exposid{val}))} and
then returns \tcode{ex\-pected<U, E>()}.
\end{itemize}
\end{itemdescr}

\indexlibrarymember{transform_error}{expected}%
\begin{itemdecl}
template<class F> constexpr auto transform_error(F&& f) &;
template<class F> constexpr auto transform_error(F&& f) const &;
\end{itemdecl}

\begin{itemdescr}
\pnum
Let \tcode{G} be \tcode{remove_cv_t<invoke_result_t<F, decltype(error())>>}.

\pnum
\constraints
\tcode{is_constructible_v<T, decltype((\exposid{val}))>} is \tcode{true}.

\pnum
\mandates
\tcode{G} is a valid template argument
for \tcode{unexpected}\iref{expected.un.general} and the declaration
\begin{codeblock}
G g(invoke(std::forward<F>(f), error()));
\end{codeblock}
is well-formed.

\pnum
\returns
If \tcode{has_value()} is \tcode{true},
\tcode{expected<T, G>(in_place, \exposid{val})}; otherwise, an \tcode{expected<T, G>}
object whose \exposid{has_val} member is \tcode{false} and \exposid{unex} member
is direct-non-list-initialized with \tcode{invoke(std::forward<F>(f), error())}.
\end{itemdescr}

\indexlibrarymember{transform_error}{expected}%
\begin{itemdecl}
template<class F> constexpr auto transform_error(F&& f) &&;
template<class F> constexpr auto transform_error(F&& f) const &&;
\end{itemdecl}

\begin{itemdescr}
\pnum
Let \tcode{G} be
\tcode{remove_cv_t<invoke_result_t<F, decltype(std::move(error()))>>}.

\pnum
\constraints
\tcode{is_constructible_v<T, decltype(std::move(\exposid{val}))>} is \tcode{true}.

\pnum
\mandates
\tcode{G} is a valid template argument
for \tcode{unexpected}\iref{expected.un.general} and the declaration
\begin{codeblock}
G g(invoke(std::forward<F>(f), std::move(error())));
\end{codeblock}
is well-formed.

\pnum
\returns
If \tcode{has_value()} is \tcode{true},
\tcode{expected<T, G>(in_place, std::move(\exposid{val}))}; otherwise, an
\tcode{expected<T, G>} object whose \exposid{has_val} member is \tcode{false}
and \exposid{unex} member is direct-non-list-initialized with
\tcode{invoke(std::forward<F>(f), std::move(error()))}.
\end{itemdescr}

\rSec3[expected.object.eq]{Equality operators}

\indexlibrarymember{operator==}{expected}%
\begin{itemdecl}
template<class T2, class E2> requires (!is_void_v<T2>)
  friend constexpr bool operator==(const expected& x, const expected<T2, E2>& y);
\end{itemdecl}

\begin{itemdescr}
\pnum
\constraints
The expressions \tcode{*x == *y} and \tcode{x.error() == y.error()}
are well-formed and their results are convertible to \tcode{bool}.

\pnum
\returns
If \tcode{x.has_value()} does not equal \tcode{y.has_value()}, \tcode{false};
otherwise if \tcode{x.has_value()} is \tcode{true}, \tcode{*x == *y};
otherwise \tcode{x.error() == y.error()}.
\end{itemdescr}

\indexlibrarymember{operator==}{expected}%
\begin{itemdecl}
template<class T2> friend constexpr bool operator==(const expected& x, const T2& v);
\end{itemdecl}

\begin{itemdescr}
\pnum
\constraints
\tcode{T2} is not a specialization of \tcode{expected}.
The expression \tcode{*x == v} is well-formed and
its result is convertible to \tcode{bool}.
\begin{note}
\tcode{T} need not be \oldconcept{EqualityComparable}.
\end{note}

\pnum
\returns
If \tcode{x.has_value()} is \tcode{true},
\tcode{*x == v};
otherwise \tcode{false}.
\end{itemdescr}

\indexlibrarymember{operator==}{expected}%
\begin{itemdecl}
template<class E2> friend constexpr bool operator==(const expected& x, const unexpected<E2>& e);
\end{itemdecl}

\begin{itemdescr}
\pnum
\constraints
The expression \tcode{x.error() == e.error()} is well-formed and
its result is convertible to \tcode{bool}.

\pnum
\returns
If \tcode{!x.has_value()} is \tcode{true},
\tcode{x.error() == e.error()};
otherwise \tcode{false}.
\end{itemdescr}

\rSec2[expected.void]{Partial specialization of \tcode{expected} for \tcode{void} types}

\rSec3[expected.void.general]{General}

\begin{codeblock}
template<class T, class E> requires is_void_v<T>
class expected<T, E> {
public:
  using @\libmember{value_type}{expected<void>}@ = T;
  using @\libmember{error_type}{expected<void>}@ = E;
  using @\libmember{unexpected_type}{expected<void>}@ = unexpected<E>;

  template<class U>
  using @\libmember{rebind}{expected<void>}@ = expected<U, error_type>;

  // \ref{expected.void.cons}, constructors
  constexpr expected() noexcept;
  constexpr expected(const expected&);
  constexpr expected(expected&&) noexcept(@\seebelow@);
  template<class U, class G>
    constexpr explicit(@\seebelow@) expected(const expected<U, G>&);
  template<class U, class G>
    constexpr explicit(@\seebelow@) expected(expected<U, G>&&);

  template<class G>
    constexpr explicit(@\seebelow@) expected(const unexpected<G>&);
  template<class G>
    constexpr explicit(@\seebelow@) expected(unexpected<G>&&);

  constexpr explicit expected(in_place_t) noexcept;
  template<class... Args>
    constexpr explicit expected(unexpect_t, Args&&...);
  template<class U, class... Args>
    constexpr explicit expected(unexpect_t, initializer_list<U>, Args&&...);

  // \ref{expected.void.dtor}, destructor
  constexpr ~expected();

  // \ref{expected.void.assign}, assignment
  constexpr expected& operator=(const expected&);
  constexpr expected& operator=(expected&&) noexcept(@\seebelow@);
  template<class G>
    constexpr expected& operator=(const unexpected<G>&);
  template<class G>
    constexpr expected& operator=(unexpected<G>&&);
  constexpr void emplace() noexcept;

  // \ref{expected.void.swap}, swap
  constexpr void swap(expected&) noexcept(@\seebelow@);
  friend constexpr void swap(expected& x, expected& y) noexcept(noexcept(x.swap(y)));

  // \ref{expected.void.obs}, observers
  constexpr explicit operator bool() const noexcept;
  constexpr bool has_value() const noexcept;
  constexpr void operator*() const noexcept;
  constexpr void value() const &;                                           // freestanding-deleted
  constexpr void value() &&;                                                // freestanding-deleted
  constexpr const E& error() const & noexcept;
  constexpr E& error() & noexcept;
  constexpr const E&& error() const && noexcept;
  constexpr E&& error() && noexcept;
  template<class G = E> constexpr E error_or(G&&) const &;
  template<class G = E> constexpr E error_or(G&&) &&;

  // \ref{expected.void.monadic}, monadic operations
  template<class F> constexpr auto and_then(F&& f) &;
  template<class F> constexpr auto and_then(F&& f) &&;
  template<class F> constexpr auto and_then(F&& f) const &;
  template<class F> constexpr auto and_then(F&& f) const &&;
  template<class F> constexpr auto or_else(F&& f) &;
  template<class F> constexpr auto or_else(F&& f) &&;
  template<class F> constexpr auto or_else(F&& f) const &;
  template<class F> constexpr auto or_else(F&& f) const &&;
  template<class F> constexpr auto transform(F&& f) &;
  template<class F> constexpr auto transform(F&& f) &&;
  template<class F> constexpr auto transform(F&& f) const &;
  template<class F> constexpr auto transform(F&& f) const &&;
  template<class F> constexpr auto transform_error(F&& f) &;
  template<class F> constexpr auto transform_error(F&& f) &&;
  template<class F> constexpr auto transform_error(F&& f) const &;
  template<class F> constexpr auto transform_error(F&& f) const &&;

  // \ref{expected.void.eq}, equality operators
  template<class T2, class E2> requires is_void_v<T2>
    friend constexpr bool operator==(const expected& x, const expected<T2, E2>& y);
  template<class E2>
    friend constexpr bool operator==(const expected&, const unexpected<E2>&);

private:
  bool @\exposid{has_val}@;         // \expos
  union {
    unexpected<E> @\exposid{unex}@; // \expos
  };
};
\end{codeblock}

\pnum
Any object of type \tcode{expected<T, E>} either
represents a value of type \tcode{T}, or
contains a value of type \tcode{E}
nested within\iref{intro.object} it.
Member \exposid{has_val} indicates whether the \tcode{expected<T, E>} object
represents a value of type \tcode{T}.

\begin{addedblock}
\pnum
When \tcode{has_value()} is \tcode{false}, the error is
\tcode{\exposid{unex}.error()}. The error is held as an \tcode{unexpected<E>},
and not as an \tcode{E}, so that \tcode{E} may be an lvalue reference type: an
\tcode{E\&} cannot be a union member, whereas \tcode{unexpected<E\&>} holds a
pointer to an external object\iref{expected.un.ref}.
\end{addedblock}

\pnum
A program that instantiates
the definition of the template \tcode{expected<T, E>} with
a type for the \tcode{E} parameter that
is not a valid template argument for \tcode{unexpected} is ill-formed.

\pnum
\added{If \tcode{E} is not an lvalue reference type, }\tcode{E} shall meet the requirements of
\oldconcept{Destructible} (\tref{cpp17.destructible}).

\rSec3[expected.void.cons]{Constructors}

\indexlibraryctor{expected<void>}%
\begin{itemdecl}
constexpr expected() noexcept;
\end{itemdecl}

\begin{itemdescr}
\pnum
\ensures
\tcode{has_value()} is \tcode{true}.
\end{itemdescr}

\indexlibraryctor{expected<void>}%
\begin{itemdecl}
constexpr expected(const expected& rhs);
\end{itemdecl}

\begin{itemdescr}
\pnum
\effects
If \tcode{rhs.has_value()} is \tcode{false},
direct-non-list-initializes \exposid{unex} with \tcode{rhs.error()}.

\pnum
\ensures
\tcode{rhs.has_value() == this->has_value()}.

\pnum
\throws
Any exception thrown by the initialization of \exposid{unex}.

\pnum
\remarks
This constructor is defined as deleted
unless \tcode{is_copy_constructible_v<E>} is \tcode{true}.

\pnum
This constructor is trivial
if \tcode{is_trivially_copy_constructible_v<E>} is \tcode{true}.
\end{itemdescr}

\indexlibraryctor{expected<void>}%
\begin{itemdecl}
constexpr expected(expected&& rhs) noexcept(is_nothrow_move_constructible_v<E>);
\end{itemdecl}

\begin{itemdescr}
\pnum
\constraints
\tcode{is_move_constructible_v<E>} is \tcode{true}.

\pnum
\effects
If \tcode{rhs.has_value()} is \tcode{false},
direct-non-list-initializes \exposid{unex} with \tcode{std::move(rhs.error())}.

\pnum
\ensures
\tcode{rhs.has_value()} is unchanged;
\tcode{rhs.has_value() == this->has_value()} is \tcode{true}.

\pnum
\throws
Any exception thrown by the initialization of \exposid{unex}.

\pnum
\remarks
This constructor is trivial
if \tcode{is_trivially_move_constructible_v<E>} is \tcode{true}.
\end{itemdescr}

\indexlibraryctor{expected<void>}%
\begin{itemdecl}
template<class U, class G>
  constexpr explicit(!is_convertible_v<const G&, E>) expected(const expected<U, G>& rhs);
template<class U, class G>
  constexpr explicit(!is_convertible_v<G, E>) expected(expected<U, G>&& rhs);
\end{itemdecl}

\begin{itemdescr}
\pnum
Let \tcode{GF} be \tcode{const G\&} for the first overload and
\tcode{G} for the second overload.

\pnum
\constraints
\begin{itemize}
\item
\tcode{is_void_v<U>} is \tcode{true}; and
\item
\tcode{is_constructible_v<E, GF>} is \tcode{true}; and
\item
\tcode{is_constructible_v<unexpected<E>, expected<U, G>\&>}
is \tcode{false}; and
\item
\tcode{is_constructible_v<unexpected<E>, expected<U, G>>}
is \tcode{false}; and
\item
\tcode{is_constructible_v<unexpected<E>, const expected<U, G>\&>}
is \tcode{false}; and
\item
\tcode{is_constructible_v<unexpected<E>, const expected<U, G>>}
is \tcode{false}.
\end{itemize}

\pnum
\effects
If \tcode{rhs.has_value()} is \tcode{false},
direct-non-list-initializes \exposid{unex}
with \tcode{std::forward<GF>(rhs.er\-ror())}.

\pnum
\ensures
\tcode{rhs.has_value()} is unchanged;
\tcode{rhs.has_value() == this->has_value()} is \tcode{true}.

\pnum
\throws
Any exception thrown by the initialization of \exposid{unex}.
\end{itemdescr}

\indexlibraryctor{expected<void>}%
\begin{itemdecl}
template<class G>
  constexpr explicit(!is_convertible_v<const G&, E>) expected(const unexpected<G>& e);
template<class G>
  constexpr explicit(!is_convertible_v<G, E>) expected(unexpected<G>&& e);
\end{itemdecl}

\begin{itemdescr}
\pnum
Let \tcode{GF} be \tcode{const G\&} for the first overload and
\tcode{G} for the second overload.

\pnum
\constraints
\tcode{is_constructible_v<E, GF>} is \tcode{true}.

\pnum
\effects
Direct-non-list-initializes \exposid{unex}
with \tcode{std::forward<GF>(e.error())}.

\pnum
\ensures
\tcode{has_value()} is \tcode{false}.

\pnum
\throws
Any exception thrown by the initialization of \exposid{unex}.
\end{itemdescr}

\indexlibraryctor{expected<void>}%
\begin{itemdecl}
constexpr explicit expected(in_place_t) noexcept;
\end{itemdecl}

\begin{itemdescr}
\pnum
\ensures
\tcode{has_value()} is \tcode{true}.
\end{itemdescr}

\indexlibraryctor{expected<void>}%
\begin{itemdecl}
template<class... Args>
  constexpr explicit expected(unexpect_t, Args&&... args);
\end{itemdecl}

\begin{itemdescr}
\pnum
\constraints
\tcode{is_constructible_v<E, Args...>} is \tcode{true}.

\pnum
\effects
Direct-non-list-initializes \exposid{unex}
with \tcode{std::forward<Args>(args)...}.

\pnum
\ensures
\tcode{has_value()} is \tcode{false}.

\pnum
\throws
Any exception thrown by the initialization of \exposid{unex}.
\end{itemdescr}

\indexlibraryctor{expected<void>}%
\begin{itemdecl}
template<class U, class... Args>
  constexpr explicit expected(unexpect_t, initializer_list<U> il, Args&&... args);
\end{itemdecl}

\begin{itemdescr}
\pnum
\constraints
\tcode{is_constructible_v<E, initializer_list<U>\&, Args...>} is \tcode{true}.

\pnum
\effects
Direct-non-list-initializes \exposid{unex}
with \tcode{il, std::forward<Args>(args)...}.

\pnum
\ensures
\tcode{has_value()} is \tcode{false}.

\pnum
\throws
Any exception thrown by the initialization of \exposid{unex}.
\end{itemdescr}

\rSec3[expected.void.dtor]{Destructor}

\indexlibrarydtor{expected<void>}%
\begin{itemdecl}
constexpr ~expected();
\end{itemdecl}

\begin{itemdescr}
\pnum
\effects
If \tcode{has_value()} is \tcode{false}, destroys \exposid{unex}.

\pnum
\remarks
If \tcode{is_trivially_destructible_v<E>} is \tcode{true},
then this destructor is a trivial destructor.
\end{itemdescr}

\rSec3[expected.void.assign]{Assignment}

\indexlibrarymember{operator=}{expected<void>}%
\begin{itemdecl}
constexpr expected& operator=(const expected& rhs);
\end{itemdecl}

\begin{itemdescr}
\pnum
\effects
\begin{itemize}
\item
If \tcode{this->has_value() \&\& rhs.has_value()} is \tcode{true}, no effects.
\item
Otherwise, if \tcode{this->has_value()} is \tcode{true},
equivalent to: \tcode{construct_at(addressof(\exposid{unex}), rhs.\exposid{unex}); \exposid{has_val} = false;}
\item
Otherwise, if \tcode{rhs.has_value()} is \tcode{true},
destroys \exposid{unex} and sets \exposid{has_val} to \tcode{true}.
\item
Otherwise, equivalent to \tcode{\exposid{unex} = rhs.\exposid{unex}}.
\end{itemize}

\pnum
\returns
\tcode{*this}.

\pnum
\remarks
This operator is defined as deleted unless
\tcode{is_copy_assignable_v<E>} is \tcode{true} and
\tcode{is_copy_constructible_v<E>} is \tcode{true}.

\pnum
This operator is trivial if
\tcode{is_trivially_copy_constructible_v<E>},
\tcode{is_trivially_copy_assigna\-ble_v<E>}, and
\tcode{is_trivially_destructible_v<E>}
are all \tcode{true}.
\end{itemdescr}

\indexlibrarymember{operator=}{expected<void>}%
\begin{itemdecl}
constexpr expected& operator=(expected&& rhs) noexcept(@\seebelow@);
\end{itemdecl}

\begin{itemdescr}
\pnum
\constraints
\tcode{is_move_constructible_v<E>} is \tcode{true} and
\tcode{is_move_assignable_v<E>} is \tcode{true}.

\pnum
\effects
\begin{itemize}
\item
If \tcode{this->has_value() \&\& rhs.has_value()} is \tcode{true}, no effects.
\item
Otherwise, if \tcode{this->has_value()} is \tcode{true}, equivalent to:
\begin{codeblock}
construct_at(addressof(@\exposid{unex}@), std::move(rhs.@\exposid{unex}@));
@\exposid{has_val}@ = false;
\end{codeblock}
\item
Otherwise, if \tcode{rhs.has_value()} is \tcode{true},
destroys \exposid{unex} and sets \exposid{has_val} to \tcode{true}.
\item
Otherwise, equivalent to \tcode{\exposid{unex} = std::move(rhs.\exposid{unex})}.
\end{itemize}

\pnum
\returns
\tcode{*this}.

\pnum
\remarks
The exception specification is equivalent to
\tcode{is_nothrow_move_constructible_v<E> \&\& is_nothrow_move_assignable_v<E>}.

\pnum
This operator is trivial if
\tcode{is_trivially_move_constructible_v<E>},
\tcode{is_trivially_move_assigna\-ble_v<E>}, and
\tcode{is_trivially_destructible_v<E>}
are all \tcode{true}.
\end{itemdescr}

\indexlibrarymember{operator=}{expected<void>}%
\begin{itemdecl}
template<class G>
  constexpr expected& operator=(const unexpected<G>& e);
template<class G>
  constexpr expected& operator=(unexpected<G>&& e);
\end{itemdecl}

\begin{itemdescr}
\pnum
Let \tcode{GF} be \tcode{const G\&} for the first overload and
\tcode{G} for the second overload.

\pnum
\constraints
\tcode{is_constructible_v<E, GF>} is \tcode{true} and
\tcode{is_assignable_v<E\&, GF>} is \tcode{true}.

\pnum
\effects
\begin{itemize}
\item
If \tcode{has_value()} is \tcode{true}, equivalent to:
\begin{codeblock}
construct_at(addressof(@\exposid{unex}@), std::forward<GF>(e.error()));
@\exposid{has_val}@ = false;
\end{codeblock}
\item
Otherwise, equivalent to:
\tcode{\exposid{unex} = unexpected<E>(std::forward<GF>(e.error()));}
\end{itemize}

\pnum
\returns
\tcode{*this}.
\end{itemdescr}

\indexlibrarymember{emplace}{expected<void>}%
\begin{itemdecl}
constexpr void emplace() noexcept;
\end{itemdecl}

\begin{itemdescr}
\pnum
\effects
If \tcode{has_value()} is \tcode{false},
destroys \exposid{unex} and sets \exposid{has_val} to \tcode{true}.
\end{itemdescr}

\rSec3[expected.void.swap]{Swap}

\indexlibrarymember{swap}{expected<void>}%
\begin{itemdecl}
constexpr void swap(expected& rhs) noexcept(@\seebelow@);
\end{itemdecl}

\begin{itemdescr}
\pnum
\constraints
\tcode{is_swappable_v<E>} is \tcode{true} and
\tcode{is_move_constructible_v<E>} is \tcode{true}.

\pnum
\effects
See \tref{expected.void.swap}.

\begin{floattable}{\tcode{swap(expected\&)} effects}{expected.void.swap}
{lx{0.35\hsize}x{0.35\hsize}}
\topline
& \chdr{\tcode{this->has_value()}} & \rhdr{\tcode{!this->has_value()}} \\ \capsep
\lhdr{\tcode{rhs.has_value()}} &
  no effects &
  calls \tcode{rhs.swap(*this)} \\
\lhdr{\tcode{!rhs.has_value()}} &
  \seebelow &
  equivalent to: \tcode{using std::swap; swap(\exposid{unex}, rhs.\exposid{unex});} \\
\end{floattable}

For the case where \tcode{rhs.has_value()} is \tcode{false} and
\tcode{this->has_value()} is \tcode{true}, equivalent to:
\begin{codeblock}
construct_at(addressof(@\exposid{unex}@), std::move(rhs.@\exposid{unex}@));
destroy_at(addressof(rhs.@\exposid{unex}@));
@\exposid{has_val}@ = false;
rhs.@\exposid{has_val}@ = true;
\end{codeblock}

\pnum
\throws
Any exception thrown by the expressions in the \Fundescx{Effects}.

\pnum
\remarks
The exception specification is equivalent to
\tcode{is_nothrow_move_constructible_v<unexpected<E>> \&\& is_nothrow_swappable_v<unexpected<E>>}.
\end{itemdescr}

\indexlibrarymember{swap}{expected<void>}%
\begin{itemdecl}
friend constexpr void swap(expected& x, expected& y) noexcept(noexcept(x.swap(y)));
\end{itemdecl}

\begin{itemdescr}
\pnum
\effects
Equivalent to \tcode{x.swap(y)}.
\end{itemdescr}

\rSec3[expected.void.obs]{Observers}

\indexlibrarymember{operator bool}{expected<void>}%
\indexlibrarymember{has_value}{expected<void>}%
\begin{itemdecl}
constexpr explicit operator bool() const noexcept;
constexpr bool has_value() const noexcept;
\end{itemdecl}

\begin{itemdescr}
\pnum
\returns
\exposid{has_val}.
\end{itemdescr}

\indexlibrarymember{operator*}{expected<void>}%
\begin{itemdecl}
constexpr void operator*() const noexcept;
\end{itemdecl}

\begin{itemdescr}
\pnum
\hardexpects
\tcode{has_value()} is \tcode{true}.
\end{itemdescr}

\indexlibrarymember{value}{expected<void>}%
\begin{itemdecl}
constexpr void value() const &;
\end{itemdecl}

\begin{itemdescr}
\pnum
\mandates
\tcode{is_copy_constructible_v<E>} is \tcode{true}.

\pnum
\throws
\tcode{bad_expected_access(error())} if \tcode{has_value()} is \tcode{false}.
\end{itemdescr}

\indexlibrarymember{value}{expected<void>}%
\begin{itemdecl}
constexpr void value() &&;
\end{itemdecl}

\begin{itemdescr}
\pnum
\mandates
\tcode{is_copy_constructible_v<E>} is \tcode{true} and
\tcode{is_move_constructible_v<E>} is \tcode{true}.

\pnum
\throws
\tcode{bad_expected_access(std::move(error()))}
if \tcode{has_value()} is \tcode{false}.
\end{itemdescr}

\indexlibrarymember{error}{expected<void>}%
\begin{itemdecl}
constexpr const E& error() const & noexcept;
constexpr E& error() & noexcept;
\end{itemdecl}

\begin{itemdescr}
\pnum
\hardexpects
\tcode{has_value()} is \tcode{false}.

\pnum
\returns
\tcode{\exposid{unex}.error()}.
\end{itemdescr}

\indexlibrarymember{error}{expected<void>}%
\begin{itemdecl}
constexpr E&& error() && noexcept;
constexpr const E&& error() const && noexcept;
\end{itemdecl}

\begin{itemdescr}
\pnum
\hardexpects
\tcode{has_value()} is \tcode{false}.

\pnum
\returns
\tcode{std::move(\exposid{unex}).error()}.
\end{itemdescr}

\indexlibrarymember{error_or}{expected<void>}%
\begin{itemdecl}
template<class G = E> constexpr E error_or(G&& e) const &;
\end{itemdecl}

\begin{itemdescr}
\pnum
\mandates
\tcode{is_copy_constructible_v<E>} is \tcode{true} and
\tcode{is_convertible_v<G, E>} is \tcode{true}.

\pnum
\returns
\tcode{std::forward<G>(e)} if \tcode{has_value()} is \tcode{true},
\tcode{error()} otherwise.
\end{itemdescr}

\indexlibrarymember{error_or}{expected<void>}%
\begin{itemdecl}
template<class G = E> constexpr E error_or(G&& e) &&;
\end{itemdecl}

\begin{itemdescr}
\pnum
\mandates
\tcode{is_move_constructible_v<E>} is \tcode{true} and
\tcode{is_convertible_v<G, E>} is \tcode{true}.

\pnum
\returns
\tcode{std::forward<G>(e)} if \tcode{has_value()} is \tcode{true},
\tcode{std::move(error())} otherwise.
\end{itemdescr}

\rSec3[expected.void.monadic]{Monadic operations}

\indexlibrarymember{and_then}{expected<void>}%
\begin{itemdecl}
template<class F> constexpr auto and_then(F&& f) &;
template<class F> constexpr auto and_then(F&& f) const &;
\end{itemdecl}

\begin{itemdescr}
\pnum
Let \tcode{U} be \tcode{remove_cvref_t<invoke_result_t<F>>}.

\pnum
\constraints
\tcode{is_constructible_v<E, decltype(error())>>} is \tcode{true}.

\pnum
\mandates
\tcode{U} is a specialization of \tcode{expected} and
\tcode{is_same_v<typename U::error_type, E>} is \tcode{true}.

\pnum
\effects
Equivalent to:
\begin{codeblock}
if (has_value())
  return invoke(std::forward<F>(f));
else
  return U(unexpect, error());
\end{codeblock}
\end{itemdescr}

\indexlibrarymember{and_then}{expected<void>}%
\begin{itemdecl}
template<class F> constexpr auto and_then(F&& f) &&;
template<class F> constexpr auto and_then(F&& f) const &&;
\end{itemdecl}

\begin{itemdescr}
\pnum
Let \tcode{U} be \tcode{remove_cvref_t<invoke_result_t<F>>}.

\pnum
\constraints
\tcode{is_constructible_v<E, decltype(std::move(error()))>} is \tcode{true}.

\pnum
\mandates
\tcode{U} is a specialization of \tcode{expected} and
\tcode{is_same_v<typename U::error_type, E>} is \tcode{true}.

\pnum
\effects
Equivalent to:
\begin{codeblock}
if (has_value())
  return invoke(std::forward<F>(f));
else
  return U(unexpect, std::move(error()));
\end{codeblock}
\end{itemdescr}

\indexlibrarymember{or_else}{expected<void>}%
\begin{itemdecl}
template<class F> constexpr auto or_else(F&& f) &;
template<class F> constexpr auto or_else(F&& f) const &;
\end{itemdecl}

\begin{itemdescr}
\pnum
Let \tcode{G} be \tcode{remove_cvref_t<invoke_result_t<F, decltype(error())>>}.

\pnum
\mandates
\tcode{G} is a specialization of \tcode{expected} and
\tcode{is_same_v<typename G::value_type, T>} is \tcode{true}.

\pnum
\effects
Equivalent to:
\begin{codeblock}
if (has_value())
  return G();
else
  return invoke(std::forward<F>(f), error());
\end{codeblock}
\end{itemdescr}

\indexlibrarymember{or_else}{expected<void>}%
\begin{itemdecl}
template<class F> constexpr auto or_else(F&& f) &&;
template<class F> constexpr auto or_else(F&& f) const &&;
\end{itemdecl}

\begin{itemdescr}
\pnum
Let \tcode{G} be
\tcode{remove_cvref_t<invoke_result_t<F, decltype(std::move(error()))>>}.

\pnum
\mandates
\tcode{G} is a specialization of \tcode{expected} and
\tcode{is_same_v<typename G::value_type, T>} is \tcode{true}.

\pnum
\effects
Equivalent to:
\begin{codeblock}
if (has_value())
  return G();
else
  return invoke(std::forward<F>(f), std::move(error()));
\end{codeblock}
\end{itemdescr}

\indexlibrarymember{transform}{expected<void>}%
\begin{itemdecl}
template<class F> constexpr auto transform(F&& f) &;
template<class F> constexpr auto transform(F&& f) const &;
\end{itemdecl}

\begin{itemdescr}
\pnum
Let \tcode{U} be \tcode{remove_cv_t<invoke_result_t<F>>}.

\pnum
\constraints
\tcode{is_constructible_v<E, decltype(error())>} is \tcode{true}.

\pnum
\mandates
\tcode{U} is a valid value type for \tcode{expected}. If \tcode{is_void_v<U>} is
\tcode{false}, the declaration
\begin{codeblock}
U u(invoke(std::forward<F>(f)));
\end{codeblock}
is well-formed.

\pnum
\effects
\begin{itemize}
\item
If \tcode{has_value()} is \tcode{false}, returns
\tcode{expected<U, E>(unexpect, error())}.
\item
Otherwise, if \tcode{is_void_v<U>} is \tcode{false}, returns an
\tcode{expected<U, E>} object whose \exposid{has_val} member is \tcode{true} and
\exposid{val} member is direct-non-list-initialized with
\tcode{invoke(std::forward<F>(f))}.
\item
Otherwise, evaluates \tcode{invoke(std::forward<F>(f))} and then returns
\tcode{expected<U, E>()}.
\end{itemize}
\end{itemdescr}

\indexlibrarymember{transform}{expected<void>}%
\begin{itemdecl}
template<class F> constexpr auto transform(F&& f) &&;
template<class F> constexpr auto transform(F&& f) const &&;
\end{itemdecl}

\begin{itemdescr}
\pnum
Let \tcode{U} be \tcode{remove_cv_t<invoke_result_t<F>>}.

\pnum
\constraints
\tcode{is_constructible_v<E, decltype(std::move(error()))>} is \tcode{true}.

\pnum
\mandates
\tcode{U} is a valid value type for \tcode{expected}. If \tcode{is_void_v<U>} is
\tcode{false}, the declaration
\begin{codeblock}
U u(invoke(std::forward<F>(f)));
\end{codeblock}
is well-formed.

\pnum
\effects
\begin{itemize}
\item
If \tcode{has_value()} is \tcode{false}, returns
\tcode{expected<U, E>(unexpect, std::move(error()))}.
\item
Otherwise, if \tcode{is_void_v<U>} is \tcode{false}, returns an
\tcode{expected<U, E>} object whose \exposid{has_val} member is \tcode{true} and
\exposid{val} member is direct-non-list-initialized with
\tcode{invoke(std::forward<F>(f))}.
\item
Otherwise, evaluates \tcode{invoke(std::forward<F>(f))} and then returns
\tcode{expected<U, E>()}.
\end{itemize}
\end{itemdescr}

\indexlibrarymember{transform_error}{expected<void>}%
\begin{itemdecl}
template<class F> constexpr auto transform_error(F&& f) &;
template<class F> constexpr auto transform_error(F&& f) const &;
\end{itemdecl}

\begin{itemdescr}
\pnum
Let \tcode{G} be \tcode{remove_cv_t<invoke_result_t<F, decltype(error())>>}.

\pnum
\mandates
\tcode{G} is a valid template argument
for \tcode{unexpected}\iref{expected.un.general} and the declaration
\begin{codeblock}
G g(invoke(std::forward<F>(f), error()));
\end{codeblock}
is well-formed.

\pnum
\returns
If \tcode{has_value()} is \tcode{true}, \tcode{expected<T, G>()}; otherwise, an
\tcode{expected<T, G>} object whose \exposid{has_val} member is \tcode{false}
and \exposid{unex} member is direct-non-list-initialized with
\tcode{invoke(std::for\-ward<F>(f), error())}.
\end{itemdescr}

\indexlibrarymember{transform_error}{expected<void>}%
\begin{itemdecl}
template<class F> constexpr auto transform_error(F&& f) &&;
template<class F> constexpr auto transform_error(F&& f) const &&;
\end{itemdecl}

\begin{itemdescr}
\pnum
Let \tcode{G} be
\tcode{remove_cv_t<invoke_result_t<F, decltype(std::move(error()))>>}.

\pnum
\mandates
\tcode{G} is a valid template argument
for \tcode{unexpected}\iref{expected.un.general} and the declaration
\begin{codeblock}
G g(invoke(std::forward<F>(f), std::move(error())));
\end{codeblock}
is well-formed.

\pnum
\returns
If \tcode{has_value()} is \tcode{true}, \tcode{expected<T, G>()}; otherwise, an
\tcode{expected<T, G>} object whose \exposid{has_val} member is \tcode{false}
and \exposid{unex} member is direct-non-list-initialized with
\tcode{invoke(std::for\-ward<F>(f), std::move(error()))}.
\end{itemdescr}

\rSec3[expected.void.eq]{Equality operators}

\indexlibrarymember{operator==}{expected<void>}%
\begin{itemdecl}
template<class T2, class E2> requires is_void_v<T2>
  friend constexpr bool operator==(const expected& x, const expected<T2, E2>& y);
\end{itemdecl}

\begin{itemdescr}
\pnum
\constraints
The expression \tcode{x.error() == y.error()} is well-formed and
its result is convertible to \tcode{bool}.

\pnum
\returns
If \tcode{x.has_value()} does not equal \tcode{y.has_value()}, \tcode{false};
otherwise if \tcode{x.has_value()} is \tcode{true}, \tcode{true};
otherwise \tcode{x.error() == y.error()}.
\end{itemdescr}

\indexlibrarymember{operator==}{expected<void>}%
\begin{itemdecl}
template<class E2>
  friend constexpr bool operator==(const expected& x, const unexpected<E2>& e);
\end{itemdecl}

\begin{itemdescr}
\pnum
\constraints
The expression \tcode{x.error() == e.error()} is well-formed and
its result is convertible to \tcode{bool}.

\pnum
\returns
If \tcode{!x.has_value()} is \tcode{true},
\tcode{x.error() == e.error()};
otherwise \tcode{false}.
\end{itemdescr}

\begin{addedblock}
\rSec2[expected.ref]{Partial specialization of \tcode{expected} for reference types}

\rSec3[expected.ref.general]{General}

\begin{codeblock}
template<class T, class E>
class expected<T&, E> {
public:
  using value_type      = T&;
  using error_type      = E;
  using unexpected_type = unexpected<E>;

  template<class U>
  using rebind = expected<U, error_type>;

  // \ref{expected.ref.cons}, constructors
  constexpr expected() = delete;
  constexpr expected(const expected&);
  constexpr expected(expected&&) noexcept(@\seebelow@);
  template<class U = T>
    constexpr explicit(@\seebelow@) expected(U&&);
  template<class U>
    constexpr expected(U&&) = delete;
  template<class U, class G>
    constexpr explicit(@\seebelow@) expected(const expected<U&, G>&);
  template<class U, class G>
    constexpr explicit(@\seebelow@) expected(expected<U&, G>&&);
  template<class G>
    constexpr explicit(@\seebelow@) expected(const unexpected<G>&);
  template<class G>
    constexpr explicit(@\seebelow@) expected(unexpected<G>&&);
  template<class... Args>
    constexpr expected(in_place_t, Args&&...) = delete;
  template<class... Args>
    constexpr explicit expected(unexpect_t, Args&&...);
  template<class U, class... Args>
    constexpr explicit expected(unexpect_t, initializer_list<U>, Args&&...);

  // \ref{expected.ref.dtor}, destructor
  constexpr ~expected();

  // \ref{expected.ref.assign}, assignment
  constexpr expected& operator=(const expected&);
  constexpr expected& operator=(expected&&) noexcept(@\seebelow@);
  template<class U = T>
    constexpr expected& operator=(U&&);
  template<class G>
    constexpr expected& operator=(const unexpected<G>&);
  template<class G>
    constexpr expected& operator=(unexpected<G>&&);
  template<class U = T>
    constexpr T& emplace(U&&) noexcept(@\seebelow@);

  // \ref{expected.ref.swap}, swap
  constexpr void swap(expected&) noexcept(@\seebelow@);
  friend constexpr void swap(expected& x, expected& y) noexcept(noexcept(x.swap(y)));

  // \ref{expected.ref.obs}, observers
  constexpr T* operator->() const noexcept;
  constexpr T& operator*() const noexcept;
  constexpr explicit operator bool() const noexcept;
  constexpr bool has_value() const noexcept;
  constexpr T& value() const &;
  constexpr T& value() &&;
  constexpr const E& error() const & noexcept;
  constexpr E& error() & noexcept;
  constexpr const E&& error() const && noexcept;
  constexpr E&& error() && noexcept;
  template<class U = remove_cv_t<T>> constexpr remove_cv_t<T> value_or(U&&) const;
  template<class G = E> constexpr E error_or(G&&) const &;

  // \ref{expected.ref.monadic}, monadic operations
  template<class F> constexpr auto and_then(F&& f) &;
  template<class F> constexpr auto and_then(F&& f) &&;
  template<class F> constexpr auto and_then(F&& f) const &;
  template<class F> constexpr auto and_then(F&& f) const &&;
  template<class F> constexpr auto or_else(F&& f) &;
  template<class F> constexpr auto or_else(F&& f) &&;
  template<class F> constexpr auto or_else(F&& f) const &;
  template<class F> constexpr auto or_else(F&& f) const &&;
  template<class F> constexpr auto transform(F&& f) &;
  template<class F> constexpr auto transform(F&& f) &&;
  template<class F> constexpr auto transform(F&& f) const &;
  template<class F> constexpr auto transform(F&& f) const &&;
  template<class F> constexpr auto transform_error(F&& f) &;
  template<class F> constexpr auto transform_error(F&& f) &&;
  template<class F> constexpr auto transform_error(F&& f) const &;
  template<class F> constexpr auto transform_error(F&& f) const &&;

  // \ref{expected.ref.eq}, equality operators
  template<class T2, class E2> requires (!is_void_v<T2>)
    friend constexpr bool operator==(const expected& x, const expected<T2, E2>& y);
  template<class T2>
    friend constexpr bool operator==(const expected&, const T2&);
  template<class E2>
    friend constexpr bool operator==(const expected&, const unexpected<E2>&);

private:
  bool @\exposid{has_val}@;           // \expos
  union {
    T* @\exposid{val}@;               // \expos
    unexpected<E> @\exposid{unex}@;   // \expos
  };
};
\end{codeblock}

\pnum
An object of type \tcode{expected<T\&, E>} either represents a reference to an
object of type \tcode{T}, or holds an error.
Member \exposid{has_val} indicates whether the object represents a reference.
When it represents a reference, member \exposid{val} points to the referenced
object, which is not owned by the \tcode{expected} object.
Otherwise, the error is \tcode{\exposid{unex}.error()}.

\pnum
The error is held as an \tcode{unexpected<E>}, and not as an \tcode{E}, so that
\tcode{E} can be an lvalue reference type: an \tcode{E\&} cannot be a union
member, whereas \tcode{unexpected<E\&>} holds a pointer to an external
object\iref{expected.un.ref}.

\pnum
A program that instantiates the definition of
\tcode{expected<T\&, E>} with a type for the \tcode{E} parameter that
is not a valid template argument for \tcode{unexpected} is ill-formed.

\pnum
\tcode{T} shall be an object type that is not an array type.
If \tcode{E} is not an lvalue reference type,
\tcode{E} shall meet the \oldconcept{Destructible} requirements
(\tref{cpp17.destructible}).

\rSec3[expected.ref.cons]{Constructors}

\indexlibraryctor{expected}%
\begin{itemdecl}
constexpr expected(const expected& rhs);
constexpr expected(expected&& rhs) noexcept(@\seebelow@);
\end{itemdecl}

\begin{itemdescr}
\pnum
\effects
If \tcode{rhs.has_value()} is \tcode{true}, initializes \exposid{val} with
\tcode{rhs.\exposid{val}}, so that \tcode{*this} and \tcode{rhs} refer to the
same object; otherwise, initializes \exposid{unex} with
\tcode{rhs.\exposid{unex}}.

\pnum
\ensures
\tcode{rhs.has_value() == this->has_value()}.

\pnum
\remarks
The exception specification of the move constructor is equivalent to
\tcode{is_nothrow_move_constructible_v<unexpected<E>>}. Each of these
constructors is trivial if the corresponding constructor of
\tcode{unexpected<E>} is trivial, and is defined as deleted unless that
constructor of \tcode{unexpected<E>} is available.
\end{itemdescr}

\indexlibraryctor{expected}%
\begin{itemdecl}
template<class U = T>
  constexpr explicit(@\seebelow@) expected(U&& u);
\end{itemdecl}

\begin{itemdescr}
\pnum
\constraints
\begin{itemize}
\item \tcode{remove_cvref_t<U>} is not \tcode{in_place_t},
\tcode{expected}, or a specialization of \tcode{unexpected};
\item \tcode{is_constructible_v<T\&, U>} is \tcode{true}; and
\item \tcode{reference_constructs_from_temporary_v<T\&, U>} is \tcode{false}.
\end{itemize}

\pnum
\effects
Let \tcode{r} be the lvalue result of \tcode{T\& r = std::forward<U>(u);}.
Initializes \exposid{val} with \tcode{addressof(r)}.

\pnum
\ensures
\tcode{has_value()} is \tcode{true}.

\pnum
\remarks
The expression inside \keyword{explicit} is equivalent to
\tcode{!is_convertible_v<U, T\&>}.
A constructor for which
\tcode{reference_constructs_from_temporary_v<T\&, U>} is \tcode{true} ---
one that would bind \tcode{T\&} to a temporary --- is defined as deleted.
\end{itemdescr}

\indexlibraryctor{expected}%
\begin{itemdecl}
template<class U, class G>
  constexpr explicit(@\seebelow@) expected(const expected<U&, G>& rhs);
template<class U, class G>
  constexpr explicit(@\seebelow@) expected(expected<U&, G>&& rhs);
\end{itemdecl}

\begin{itemdescr}
\pnum
\constraints
\begin{itemize}
\item \tcode{is_constructible_v<T\&, U\&>} is \tcode{true};
\item \tcode{reference_constructs_from_temporary_v<T\&, U\&>} is \tcode{false};
and
\item if \tcode{E} is an lvalue reference type, \tcode{G} is an lvalue
reference type and \tcode{is_convertible_v<G, E>} is \tcode{true}; otherwise
\tcode{is_constructible_v<E, const G\&>} (for the first overload) or
\tcode{is_constructible_v<E, G>} (for the second) is \tcode{true}.
\end{itemize}

\pnum
\effects
If \tcode{rhs.has_value()} is \tcode{true}, initializes \exposid{val} with
\tcode{addressof(*rhs)}, so that \tcode{*this} refers to the object referred to
by \tcode{rhs}; otherwise, initializes \exposid{unex} with the error of
\tcode{rhs}. No object referred to by \tcode{rhs} is moved from.

\pnum
\remarks
The expression inside \keyword{explicit} is equivalent to
\tcode{!is_convertible_v<U\&, T\&> || !is_convertible_v<G, E>}.
\end{itemdescr}

\indexlibraryctor{expected}%
\begin{itemdecl}
template<class G>
  constexpr explicit(@\seebelow@) expected(const unexpected<G>& e);
template<class G>
  constexpr explicit(@\seebelow@) expected(unexpected<G>&& e);
\end{itemdecl}

\begin{itemdescr}
\pnum
\constraints
If \tcode{E} is an lvalue reference type, \tcode{G} is an lvalue reference
type, \tcode{is_convertible_v<G, E>} is \tcode{true}, and
\tcode{reference_constructs_from_temporary_v<E, G>} is \tcode{false};
otherwise \tcode{is_constructible_v<E, const G\&>} (for the first overload) or
\tcode{is_constructible_v<E, G>} (for the second) is \tcode{true}.

\pnum
\effects
Initializes \exposid{unex} with the error of \tcode{e}.
\tcode{has_value()} is \tcode{false}.

\pnum
\remarks
When \tcode{E} is an lvalue reference type and \tcode{G} is not a reference
type, these constructors are defined as deleted: the referent would be the
value owned by \tcode{e}, and binding \tcode{E} to it would dangle once
\tcode{e} is destroyed.
\end{itemdescr}

\indexlibraryctor{expected}%
\begin{itemdecl}
template<class... Args>
  constexpr explicit expected(unexpect_t, Args&&... args);
template<class U, class... Args>
  constexpr explicit expected(unexpect_t, initializer_list<U> il, Args&&... args);
\end{itemdecl}

\begin{itemdescr}
\pnum
\constraints
\tcode{is_constructible_v<E, Args...>} (respectively,
\tcode{is_constructible_v<E, initializer_list<U>\&, Args...>}) is \tcode{true}.
If \tcode{E} is an lvalue reference type, the initialization shall not bind
\tcode{E} to a temporary.

\pnum
\effects
Direct-non-list-initializes \exposid{unex} with \tcode{in_place} and
\tcode{std::forward<Args>(args)...} (respectively, \tcode{in_place},
\tcode{il}, and \tcode{std::forward<Args>(args)...}), so that the error
\tcode{\exposid{unex}.error()} is constructed from those arguments.
\tcode{has_value()} is \tcode{false}.
\end{itemdescr}

\rSec3[expected.ref.dtor]{Destructor}

\indexlibrarydtor{expected}%
\begin{itemdecl}
constexpr ~expected();
\end{itemdecl}

\begin{itemdescr}
\pnum
\effects
If \tcode{has_value()} is \tcode{false}, destroys \exposid{unex}.

\pnum
\remarks
This destructor is trivial if \tcode{unexpected<E>} is trivially
destructible. \tcode{T} is not destroyed; \tcode{*this} never owns the object
it refers to.
\end{itemdescr}

\rSec3[expected.ref.assign]{Assignment}

\pnum
Assignment rebinds: assigning to an \tcode{expected<T\&, E>} that holds a value
changes which object it refers to. It never assigns through to the referent.

\indexlibrarymember{operator=}{expected}%
\begin{itemdecl}
constexpr expected& operator=(const expected& rhs);
constexpr expected& operator=(expected&& rhs) noexcept(@\seebelow@);
\end{itemdecl}

\begin{itemdescr}
\pnum
\effects
If \tcode{rhs.has_value()} is \tcode{true}: if \tcode{has_value()} is
\tcode{true}, assigns \tcode{rhs.\exposid{val}} to \exposid{val}; otherwise
destroys \exposid{unex} and initializes \exposid{val} with
\tcode{rhs.\exposid{val}}. If \tcode{rhs.has_value()} is \tcode{false}, the
error of \tcode{rhs} is assigned to or used to initialize \exposid{unex}, as
for the primary template (\ref{expected.object.assign}). In every case
\tcode{*this} comes to refer to the object \tcode{rhs} refers to, or to hold the
error of \tcode{rhs}; no referenced object is assigned through.

\pnum
\returns
\tcode{*this}.
\end{itemdescr}

\indexlibrarymember{operator=}{expected}%
\begin{itemdecl}
template<class U = T>
  constexpr expected& operator=(U&& u);
\end{itemdecl}

\begin{itemdescr}
\pnum
\constraints
\tcode{remove_cvref_t<U>} is not \tcode{expected} or a specialization of
\tcode{unexpected}, \tcode{is_constructible_v<T\&, U>} is \tcode{true}, and
\tcode{reference_constructs_from_temporary_v<T\&, U>} is \tcode{false}.

\pnum
\effects
Let \tcode{r} be the lvalue result of \tcode{T\& r = std::forward<U>(u);}.
If \tcode{has_value()} is \tcode{true}, assigns \tcode{addressof(r)} to
\exposid{val}. Otherwise, destroys \exposid{unex}, initializes \exposid{val}
with \tcode{addressof(r)}, and sets \exposid{has_val} to \tcode{true}; if
binding \tcode{r} throws, \tcode{*this} is left unchanged.

\pnum
\returns
\tcode{*this}.
\end{itemdescr}

\indexlibrarymember{operator=}{expected}%
\begin{itemdecl}
template<class G>
  constexpr expected& operator=(const unexpected<G>& e);
template<class G>
  constexpr expected& operator=(unexpected<G>&& e);
\end{itemdecl}

\begin{itemdescr}
\pnum
\constraints
If \tcode{E} is an lvalue reference type, \tcode{G} is an lvalue reference
type, \tcode{is_convertible_v<G, E>} is \tcode{true}, and
\tcode{reference_constructs_from_temporary_v<E, G>} is \tcode{false};
otherwise the corresponding requirements of the primary template hold.

\pnum
\effects
Makes \tcode{*this} hold the error of \tcode{e}, reinitializing \exposid{unex}
from \tcode{e} rather than assigning through it. When \tcode{E} is a reference
type, \tcode{\exposid{unex}.error()} thereafter refers to the same object as
\tcode{e.error()}; the previously referenced object, if any, is not modified.

\pnum
\returns
\tcode{*this}.

\pnum
\remarks
When \tcode{E} is an lvalue reference type and \tcode{G} is not a reference
type, these overloads are defined as deleted.
\end{itemdescr}

\indexlibrarymember{emplace}{expected}%
\begin{itemdecl}
template<class U = T>
  constexpr T& emplace(U&& u) noexcept(@\seebelow@);
\end{itemdecl}

\begin{itemdescr}
\pnum
\constraints
\tcode{is_constructible_v<T\&, U>} is \tcode{true} and
\tcode{reference_constructs_from_temporary_v<T\&, U>} is \tcode{false}.

\pnum
\effects
Rebinds \tcode{*this} to refer to the object bound by
\tcode{T\& r = std::forward<U>(u);}: if \tcode{has_value()} is \tcode{false},
destroys \exposid{unex} first. Sets \exposid{val} to \tcode{addressof(r)} and
\exposid{has_val} to \tcode{true}.

\pnum
\returns
\tcode{*\exposid{val}}.
\end{itemdescr}

\rSec3[expected.ref.swap]{Swap}

\indexlibrarymember{swap}{expected}%
\begin{itemdecl}
constexpr void swap(expected& rhs) noexcept(@\seebelow@);
\end{itemdecl}

\begin{itemdescr}
\pnum
\effects
Exchanges the states of \tcode{*this} and \tcode{rhs}. When both hold values,
exchanges \exposid{val} and \tcode{rhs.\exposid{val}} --- the referenced
objects are not swapped. Otherwise behaves as the primary template's
\tcode{swap} does for the error.

\pnum
\remarks
The exception specification is equivalent to
\tcode{is_nothrow_move_constructible_v<unexpected<E>> \&\&}
\tcode{(is_reference_v<E> || is_nothrow_swappable_v<E>)}.
\end{itemdescr}

\indexlibrarymember{swap}{expected}%
\begin{itemdecl}
friend constexpr void swap(expected& x, expected& y) noexcept(noexcept(x.swap(y)));
\end{itemdecl}

\begin{itemdescr}
\pnum
\effects
Equivalent to \tcode{x.swap(y)}.
\end{itemdescr}

\rSec3[expected.ref.obs]{Observers}

\indexlibrarymember{operator->}{expected}%
\indexlibrarymember{operator*}{expected}%
\begin{itemdecl}
constexpr T* operator->() const noexcept;
constexpr T& operator*() const noexcept;
\end{itemdecl}

\begin{itemdescr}
\pnum
\expects
\tcode{has_value()} is \tcode{true}.

\pnum
\returns
\tcode{\exposid{val}} and \tcode{*\exposid{val}} respectively.

\pnum
\remarks
These are \tcode{const} member functions that return a non-\tcode{const}
\tcode{T*} and \tcode{T\&}; the constness of \tcode{*this} does not propagate
to the referenced object. For deep \tcode{const}, use
\tcode{expected<const T\&, E>}.
\end{itemdescr}

\indexlibrarymember{has_value}{expected}%
\indexlibrarymember{operator bool}{expected}%
\begin{itemdecl}
constexpr explicit operator bool() const noexcept;
constexpr bool has_value() const noexcept;
\end{itemdecl}

\begin{itemdescr}
\pnum
\returns
\tcode{\exposid{has_val}}.
\end{itemdescr}

\indexlibrarymember{value}{expected}%
\begin{itemdecl}
constexpr T& value() const &;
constexpr T& value() &&;
\end{itemdecl}

\begin{itemdescr}
\pnum
\returns
\tcode{*\exposid{val}} if \tcode{has_value()} is \tcode{true}.

\pnum
\throws
\tcode{bad_expected_access(as_const(error()))} (respectively,
\tcode{bad_expected_access(std::move(error()))}) if \tcode{has_value()} is
\tcode{false}. Both overloads yield \tcode{T\&}; the value category of
\tcode{*this} does not affect the returned reference.
\end{itemdescr}

\indexlibrarymember{error}{expected}%
\begin{itemdecl}
constexpr const E& error() const & noexcept;
constexpr E& error() & noexcept;
constexpr const E&& error() const && noexcept;
constexpr E&& error() && noexcept;
\end{itemdecl}

\begin{itemdescr}
\pnum
\expects
\tcode{has_value()} is \tcode{false}.

\pnum
\returns
\tcode{\exposid{unex}.error()}, with the \tcode{const} and reference
qualification of the selected overload.
\end{itemdescr}

\indexlibrarymember{value_or}{expected}%
\begin{itemdecl}
template<class U = remove_cv_t<T>>
  constexpr remove_cv_t<T> value_or(U&& v) const;
\end{itemdecl}

\begin{itemdescr}
\pnum
\mandates
\tcode{is_convertible_v<T\&, remove_cv_t<T>>} and
\tcode{is_convertible_v<U, remove_cv_t<T>>} are \tcode{true}.

\pnum
\returns
\tcode{has_value() ? static_cast<remove_cv_t<T>>(*\exposid{val}) :
static_cast<remove_cv_t<T>>(std::forward<U>(v))}. The result is an object,
never a reference.
\end{itemdescr}

\indexlibrarymember{error_or}{expected}%
\begin{itemdecl}
template<class G = E>
  constexpr E error_or(G&& e) const &;
\end{itemdecl}

\begin{itemdescr}
\pnum
\returns
\tcode{has_value() ? std::forward<G>(e) : error()}, as for the primary
template.
\end{itemdescr}

\rSec3[expected.ref.monadic]{Monadic operations}

\pnum
The member templates \tcode{and_then}, \tcode{or_else}, \tcode{transform}, and
\tcode{transform_error} behave as specified for the primary template
(\ref{expected.object.monadic}), with one difference: the value is passed to
the callable as \tcode{T\&} for every ref-qualification of \tcode{*this}. An
rvalue \tcode{expected<T\&, E>} does not pass its referent as an rvalue; the
object referred to is never moved from by these operations.

\rSec3[expected.ref.eq]{Equality operators}

\pnum
The equality operators behave as specified for the primary template
(\ref{expected.object.eq}), comparing referents through \tcode{operator*} and
errors through \tcode{error()}.
\end{addedblock}

%%% Local Variables:
%%% mode: LaTeX
%%% TeX-master: t
%%% End:
